% Complete reviewed Unit 14; editorial master: chapters/ch14/student.md
% Markdown SHA-256: 0f83585b45a348df5ab346d05537f1537d61fe611ee331c293eb6d09f932458d
% One-use initial assembly followed by complete content and page review.
\newcommand{\runningtitle}{Unit 14. Word formation and lexical families}
% Layout inherited from the reviewed previous course; set a running title per unit.
% This is a small hand-maintained TeX layout, not a conversion pipeline.
\RequirePackage{fix-cm}
\documentclass[11pt,twoside]{article}
\usepackage{fontspec}
\setmainfont{Cambria}[
  Renderer=HarfBuzz,
  BoldFont=cambriab.ttf,
  ItalicFont=cambriai.ttf,
  BoldItalicFont=cambriaz.ttf
]
\setmonofont{lmmono10-regular.otf}[BoldFont=lmmonolt10-bold.otf]
\usepackage[paperwidth=165mm,paperheight=241mm,inner=21mm,outer=19mm,
  top=22mm,bottom=20mm,headheight=14pt,headsep=7mm,footskip=10mm]{geometry}
\usepackage{xcolor,array,booktabs,longtable,tabularx}
\usepackage{enumitem,fancyhdr}
\usepackage[unicode,hidelinks]{hyperref}
\definecolor{ink}{HTML}{000000}
\definecolor{muted}{HTML}{4C5558}
\setlength{\parindent}{0pt}
\setlength{\parskip}{3.5pt}
\AtBeginDocument{\fontsize{11}{14.3}\selectfont}
\setlength{\tabcolsep}{4pt}
\renewcommand{\arraystretch}{1.10}
\setlist[itemize]{leftmargin=1.3em,label=\textbullet,itemsep=2pt,topsep=3pt,parsep=0pt}
\setlist[enumerate]{leftmargin=1.8em,itemsep=3pt,topsep=4pt,parsep=0pt}
\widowpenalty=10000
\clubpenalty=10000
\emergencystretch=1em
\raggedbottom
\pagestyle{fancy}
\fancyhf{}
\providecommand{\runningtitle}{Old English}
\fancyhead[LE,RO]{\small\itshape \runningtitle}
\fancyhead[LO,RE]{\small Old English}
\fancyfoot[LE,RO]{\small\thepage}
\renewcommand{\headrulewidth}{0.3pt}
\newcommand{\booktitle}[2]{
  {\fontsize{9}{11}\selectfont\color{muted}#1\par}
  \vspace{3pt}
  {\fontsize{20}{24}\selectfont\color{ink}#2\par}
  \vspace{6pt}
}
\newcommand{\source}[1]{\par{\footnotesize\color{muted}#1\par}}
\newcommand{\readingtext}[1]{
  \par\addvspace{5pt}
  \begingroup\leftskip=1.5em
    \fontsize{12}{16}\selectfont #1\par
  \endgroup\addvspace{5pt}
}
\makeatletter
\renewcommand{\section}{\@startsection{section}{1}{\z@}
  {-2.4ex \@plus -0.6ex \@minus -0.2ex}{0.8ex \@plus 0.2ex}
  {\normalfont\fontsize{12}{15}\selectfont\bfseries\color{ink}}}
\renewcommand{\subsection}{\@startsection{subsection}{2}{\z@}
  {-1.8ex \@plus -0.5ex \@minus -0.2ex}{0.5ex \@plus 0.2ex}
  {\normalfont\fontsize{11}{14}\selectfont\bfseries\color{ink}}}
\makeatother

\hyphenpenalty=10000
\exhyphenpenalty=10000
\emergencystretch=3em
\newcommand{\itempara}[2]{\par\addvspace{3pt}\noindent\hangindent=1.8em\hangafter=1\makebox[1.8em][l]{#1.}#2\par}
\newcommand{\reserve}[1]{\par\penalty-100\begingroup\dimen0=\pagegoal\advance\dimen0 by-\pagetotal\ifdim\dimen0<#1\newpage\fi\endgroup}
\begin{document}
\addfontfeatures{Ligatures=TeXOff}
\raggedright
\booktitle{UNIT 14}{Word formation and lexical families}

% mdparagraph:P001:begin
A familiar part can help you understand a word. It does not settle the whole word’s meaning. This unit distinguishes compounds and derived words from their inflected forms, then applies that distinction to the complete Cook exercises and a connected argument about the dangers of high station.
% mdparagraph:P001:end

% mdparagraph:P002:begin
Sections 1–5 supply the forms, meanings and uses needed by the practice. The four translated examples and Exercises F, A, C and S are new teaching material. Cook VIII.26–27 retain their original prompts. The reading preserves the selected Old English; its modern orientation and helps are separate additions.
% mdparagraph:P002:end

% mdparagraph:P003:begin
\textbf{Scope and notation.} Productive forms may be used in the senses and constructions taught here or earlier. A table marked “plural reference” or “singular reference” supplies that number for recognition only. Infinitive and principal-part citations identify a verb; they do not by themselves assign a new syntactic use. Supplied help stays visible during its task and is not tested by unaided recall. In verbal tables, 1, 2 and 3 denote person. Preserve vowel length and the distinction between e and ę.
% mdparagraph:P003:end

\reserve{7\baselineskip}
% mdparagraph:P004:begin
\section*{1. Formation and inflection}
% mdparagraph:P004:end

% teaching: NEW-U14-GRAM-U14-DERIVATION-VERSUS-INFLECTION
% teaching: NEW-U14-GRAM-U14-COMPOUND-ANALYSIS
% mdparagraph:P005:begin
\textbf{Derivation} forms a lexical word; \textbf{inflection} gives that word a form required in a sentence. Rihtwīs, “righteous”, and rihtwīsnes, “righteousness”, are different words and different parts of speech. The -e of rihtwīsnesse instead marks an inflected form of the noun. An ending can serve more than one case: the clause, agreement and governing word determine which one is present.
% mdparagraph:P005:end

% mdparagraph:P006:begin
A \textbf{compound} joins lexical components. In the selected noun compounds the final noun is the head: it identifies the kind of thing named and supplies the inflectional pattern. Goldhord combines gold and the supplied citation hord, “treasure; hoard”. The whole noun means “treasure; hoard of gold”; the ending is added after the complete stem goldhord-. A phrase containing an adjective and a noun instead inflects the adjective for agreement as well as the noun for case and number.
% mdparagraph:P006:end

% mdparagraph:P007:begin
The first component can specify material, source or a quality. Heofonfȳre has heofon + fȳr; hēaclifum has hēah + the supplied clif, “cliff”, with hēa- as the combining form. The local helps give these source forms’ meanings and cases. Do not turn that particular hēah/\allowbreak{}hēa- relationship into a rule for dropping h elsewhere. Learn the attested whole-word sense: gōdwębb means “purple cloth”, and wyrtgemang “spice; mixture of spices”. English component substitution will not reliably recover such meanings.\textsuperscript{1}
% mdparagraph:P007:end

% teaching: NEW-U14-GRAM-U14-DERIVED-NOUNS
\reserve{15\baselineskip}
% mdparagraph:P008:begin
A suffix follows the family stem. The selected noun suffixes include -ung for an activity, -scipe for a state or relationship, -nes for an abstract quality, and -dōm for a state, position or authority. Agent nouns in -end name a doer. These are aids to analysing attested words, not instructions to invent a noun from any base. The derived noun has its own gender, meaning and declension.
% mdparagraph:P008:end

% mdtable:T01:begin
\par\begingroup
\fontsize{10.5}{13.2}\selectfont
\setlength{\LTleft}{0pt}\setlength{\LTright}{\fill}
\noindent
\begin{tabular}{@{}>{\raggedright\arraybackslash}p{22.05mm}>{\raggedright\arraybackslash}p{52.92mm}>{\raggedright\arraybackslash}p{43.41mm}@{}}
\toprule
\textbf{Family} & \textbf{Parts of speech and learned meanings} & \textbf{Formation} \\
\midrule
leornian → leornung & verb “learn” → feminine noun “learning; study” & leorn- + -ung \\
gefēra → gefērscipe & masculine noun “companion” → masculine noun “companionship” & gefēr- + -scipe; the citation -a is not retained \\
rihtwīs → rihtwīsnes & adjective “righteous” → feminine noun “righteousness” & rihtwīs- + -nes \\
īdel → īdelnes & adjective “empty; void” → feminine noun “idleness” & īdel- + -nes; learn the noun’s established meaning \\
\bottomrule
\end{tabular}
\par\endgroup

% mdtable:T01:end
% mdparagraph:P009:begin
Leornian is supplied here as a dictionary citation of the earlier verb. Īdelnes demonstrates the limit of a family resemblance: the previously learned noun means “idleness”, not automatically “physical emptiness”. In nouns ending in -nes, s is doubled before endings: -nesse, -nessa, -nessum. That stem spelling belongs to the noun’s inflection.
% mdparagraph:P009:end

% teaching: NEW-U14-GRAM-U14-DERIVED-ADJECTIVES
% teaching: NEW-U14-GRAM-U14-ADVERB-FAMILIES
\reserve{12\baselineskip}
% mdparagraph:P010:begin
The adjective-forming -lic gives heofonlic, “heavenly”, from heofon, “heaven”. The adjective then takes the ordinary strong or weak endings. In an inflected -lic adjective, separate the suffix from the following agreement ending. The source adjective eorðlican is analysed with its supplied case, number and gender below.
% mdparagraph:P010:end

% mdtable:T02:begin
\par\begingroup
\fontsize{10.5}{13.2}\selectfont
\setlength{\LTleft}{0pt}\setlength{\LTright}{\fill}
\noindent
\begin{tabular}{@{}>{\raggedright\arraybackslash}p{41.26mm}>{\raggedright\arraybackslash}p{38.56mm}>{\raggedright\arraybackslash}p{38.56mm}@{}}
\toprule
\textbf{Citation base} & \textbf{Adverb} & \textbf{Meaning} \\
\midrule
swīð, “strong” & swīðe & very \\
sōð, “true” & sōðlice & truly \\
\bottomrule
\end{tabular}
\par\endgroup

% mdtable:T02:end
% mdparagraph:P011:begin
Swīð and sōð are supplied adjective citations, not new full adjective paradigms. The attested adverbs use -e or -lice and do not agree with a noun. A final e is not enough to identify an adverb: it also occurs in many noun and adjective forms. Use the whole word and its clause function.\textsuperscript{2}
% mdparagraph:P011:end

% teaching: NEW-U14-GRAM-U14-PREFIX-LEXICAL-LIMITS
\reserve{10\baselineskip}
% mdparagraph:P012:begin
A prefix precedes the family stem. In rihtwīsnes /\allowbreak{} unrihtwīsnes, “righteousness /\allowbreak{} unrighteousness”, un- marks opposition. Compare these supplied verb citations with the familiar finite verbs:
% mdparagraph:P012:end

% mdtable:T03:begin
\par\begingroup
\fontsize{10.5}{13.2}\selectfont
\setlength{\LTleft}{0pt}\setlength{\LTright}{\fill}
\noindent
\begin{tabular}{@{}>{\raggedright\arraybackslash}p{20.69mm}>{\raggedright\arraybackslash}p{37.43mm}>{\raggedright\arraybackslash}p{60.26mm}@{}}
\toprule
\textbf{Base} & \textbf{Prefixed word} & \textbf{Learned relation} \\
\midrule
berstan, “burst” & tōberstan, “break up; go to pieces” & tō- has an asunder sense here \\
weorðan, “become” & forweorðan, “perish” & the prefixed whole word has a destructive-result meaning \\
\bottomrule
\end{tabular}
\par\endgroup

% mdtable:T03:end
% mdparagraph:P013:begin
Each whole verb keeps its own meaning, stem changes and government. The prefix for- does not simply equal English “for”, and tō- does not always mean the same thing. Likewise, ge- can belong to a verb’s lexical form or occur in a participle; neither its addition nor its removal is automatic. Use the selected forms below.
% mdparagraph:P013:end

% teaching: NEW-U14-GRAM-U14-NOMINALIZED-EXPRESSIONS
% mdparagraph:P014:begin
An adjective can stand for people without an expressed noun: the earlier ðā rihtwīsan means “the righteous people”. It remains an adjective, with weak inflection after the article and case determined by the clause. The new neuter noun gōd, used in the plural for “goods; property”, has its own lexical entry and noun paradigm. Endings can coincide across these categories; the absence of a following noun does not decide which category a word belongs to.\textsuperscript{3}
% mdparagraph:P014:end

\reserve{7\baselineskip}
% mdparagraph:P015:begin
\section*{2. Nouns}
% mdparagraph:P015:end

% mdparagraph:P016:begin
The following tables give every singular and plural case. Gender and the ordinary meaning are stated with each word. The singular/\allowbreak{}plural reference labels limit lexical use; they do not omit cells. Otherwise both numbers are productive. The strong nouns here use genitive plural -a and dative plural -um; the weak nouns use genitive plural -ena and dative plural -um.\textsuperscript{4}
% mdparagraph:P016:end

\reserve{7\baselineskip}
% mdparagraph:P017:begin
\subsection*{Strong masculine nouns}
% mdparagraph:P017:end

% mdparagraph:P018:begin
The nominative and accusative singular have no ending. Add -es for genitive singular and -e for dative singular; nominative and accusative plural take -as. The other plural endings are -a and -um.
% mdparagraph:P018:end

% teaching: NEW-U14-VOC-FLOD
% teaching: NEW-U14-NOM-FLOD
\reserve{11\baselineskip}
% mdparagraph:P019:begin
\textbf{flōd} — masculine; flood; stream of water. Stem \textbf{flōd-}.
% mdparagraph:P019:end

% table: NEW-U14-NOM-FLOD
% mdtable:T04:begin
\par\begingroup
\fontsize{10.5}{13.2}\selectfont
\setlength{\LTleft}{0pt}\setlength{\LTright}{\fill}
\noindent
\begin{tabular}{@{}>{\raggedright\arraybackslash}p{39.46mm}>{\raggedright\arraybackslash}p{39.46mm}>{\raggedright\arraybackslash}p{39.46mm}@{}}
\toprule
\textbf{Case} & \textbf{Singular} & \textbf{Plural} \\
\midrule
Nominative & flōd & flōdas \\
Genitive & flōdes & flōda \\
Dative & flōde & flōdum \\
Accusative & flōd & flōdas \\
\bottomrule
\end{tabular}
\par\endgroup

% mdtable:T04:end
% teaching: NEW-U14-VOC-MUNT
% teaching: NEW-U14-NOM-MUNT
\reserve{11\baselineskip}
% mdparagraph:P020:begin
\textbf{munt} — masculine; mountain. Stem \textbf{munt-}.
% mdparagraph:P020:end

% table: NEW-U14-NOM-MUNT
% mdtable:T05:begin
\par\begingroup
\fontsize{10.5}{13.2}\selectfont
\setlength{\LTleft}{0pt}\setlength{\LTright}{\fill}
\noindent
\begin{tabular}{@{}>{\raggedright\arraybackslash}p{39.46mm}>{\raggedright\arraybackslash}p{39.46mm}>{\raggedright\arraybackslash}p{39.46mm}@{}}
\toprule
\textbf{Case} & \textbf{Singular} & \textbf{Plural} \\
\midrule
Nominative & munt & muntas \\
Genitive & muntes & munta \\
Dative & munte & muntum \\
Accusative & munt & muntas \\
\bottomrule
\end{tabular}
\par\endgroup

% mdtable:T05:end
% teaching: NEW-U14-VOC-HUND-DOG
% teaching: NEW-U14-NOM-HUND-DOG
\reserve{11\baselineskip}
% mdparagraph:P021:begin
\textbf{hund} — masculine; dog. Stem \textbf{hund-}.
% mdparagraph:P021:end

% table: NEW-U14-NOM-HUND-DOG
% mdtable:T06:begin
\par\begingroup
\fontsize{10.5}{13.2}\selectfont
\setlength{\LTleft}{0pt}\setlength{\LTright}{\fill}
\noindent
\begin{tabular}{@{}>{\raggedright\arraybackslash}p{39.46mm}>{\raggedright\arraybackslash}p{39.46mm}>{\raggedright\arraybackslash}p{39.46mm}@{}}
\toprule
\textbf{Case} & \textbf{Singular} & \textbf{Plural} \\
\midrule
Nominative & hund & hundas \\
Genitive & hundes & hunda \\
Dative & hunde & hundum \\
Accusative & hund & hundas \\
\bottomrule
\end{tabular}
\par\endgroup

% mdtable:T06:end
% mdparagraph:P022:begin
This hund is the noun “dog”, distinct from the earlier numeral hund.
% mdparagraph:P022:end

% teaching: NEW-U14-VOC-MIDDANGEARD
% teaching: NEW-U14-NOM-MIDDANGEARD
\reserve{11\baselineskip}
% mdparagraph:P023:begin
\textbf{middangeard} — masculine; world; earth. Stem \textbf{middangeard-}. Plural reference; singular productive.
% mdparagraph:P023:end

% table: NEW-U14-NOM-MIDDANGEARD
% mdtable:T07:begin
\par\begingroup
\fontsize{10.5}{13.2}\selectfont
\setlength{\LTleft}{0pt}\setlength{\LTright}{\fill}
\noindent
\begin{tabular}{@{}>{\raggedright\arraybackslash}p{32.88mm}>{\raggedright\arraybackslash}p{42.75mm}>{\raggedright\arraybackslash}p{42.75mm}@{}}
\toprule
\textbf{Case} & \textbf{Singular} & \textbf{Plural} \\
\midrule
Nominative & middangeard & middangeardas \\
Genitive & middangeardes & middangearda \\
Dative & middangearde & middangeardum \\
Accusative & middangeard & middangeardas \\
\bottomrule
\end{tabular}
\par\endgroup

% mdtable:T07:end
\reserve{7\baselineskip}
% mdparagraph:P024:begin
\subsection*{Strong neuter nouns}
% mdparagraph:P024:end

% mdparagraph:P025:begin
In these heavy-stem neuters the nominative and accusative are alike in both numbers and have no ending. Genitive and dative singular take -es and -e; genitive and dative plural take -a and -um. Add the ending to the whole compound stem.
% mdparagraph:P025:end

% teaching: NEW-U14-VOC-DUST
% teaching: NEW-U14-NOM-DUST
\reserve{11\baselineskip}
% mdparagraph:P026:begin
\textbf{dūst} — neuter; dust. Stem \textbf{dūst-}. Plural reference; singular productive.
% mdparagraph:P026:end

% table: NEW-U14-NOM-DUST
% mdtable:T08:begin
\par\begingroup
\fontsize{10.5}{13.2}\selectfont
\setlength{\LTleft}{0pt}\setlength{\LTright}{\fill}
\noindent
\begin{tabular}{@{}>{\raggedright\arraybackslash}p{39.46mm}>{\raggedright\arraybackslash}p{39.46mm}>{\raggedright\arraybackslash}p{39.46mm}@{}}
\toprule
\textbf{Case} & \textbf{Singular} & \textbf{Plural} \\
\midrule
Nominative & dūst & dūst \\
Genitive & dūstes & dūsta \\
Dative & dūste & dūstum \\
Accusative & dūst & dūst \\
\bottomrule
\end{tabular}
\par\endgroup

% mdtable:T08:end
% teaching: NEW-U14-VOC-WYRTGEMANG
% teaching: NEW-U14-NOM-WYRTGEMANG
\reserve{11\baselineskip}
% mdparagraph:P027:begin
\textbf{wyrtgemang} — neuter; spice; mixture of spices. Stem \textbf{wyrtgemang-}.
% mdparagraph:P027:end

% table: NEW-U14-NOM-WYRTGEMANG
% mdtable:T09:begin
\par\begingroup
\fontsize{10.5}{13.2}\selectfont
\setlength{\LTleft}{0pt}\setlength{\LTright}{\fill}
\noindent
\begin{tabular}{@{}>{\raggedright\arraybackslash}p{34.82mm}>{\raggedright\arraybackslash}p{41.78mm}>{\raggedright\arraybackslash}p{41.78mm}@{}}
\toprule
\textbf{Case} & \textbf{Singular} & \textbf{Plural} \\
\midrule
Nominative & wyrtgemang & wyrtgemang \\
Genitive & wyrtgemanges & wyrtgemanga \\
Dative & wyrtgemange & wyrtgemangum \\
Accusative & wyrtgemang & wyrtgemang \\
\bottomrule
\end{tabular}
\par\endgroup

% mdtable:T09:end
% teaching: NEW-U14-VOC-GOLDHORD
% teaching: NEW-U14-NOM-GOLDHORD
\reserve{11\baselineskip}
% mdparagraph:P028:begin
\textbf{goldhord} — neuter; treasure; hoard of gold. Stem \textbf{goldhord-}.
% mdparagraph:P028:end

% table: NEW-U14-NOM-GOLDHORD
% mdtable:T10:begin
\par\begingroup
\fontsize{10.5}{13.2}\selectfont
\setlength{\LTleft}{0pt}\setlength{\LTright}{\fill}
\noindent
\begin{tabular}{@{}>{\raggedright\arraybackslash}p{39.46mm}>{\raggedright\arraybackslash}p{39.46mm}>{\raggedright\arraybackslash}p{39.46mm}@{}}
\toprule
\textbf{Case} & \textbf{Singular} & \textbf{Plural} \\
\midrule
Nominative & goldhord & goldhord \\
Genitive & goldhordes & goldhorda \\
Dative & goldhorde & goldhordum \\
Accusative & goldhord & goldhord \\
\bottomrule
\end{tabular}
\par\endgroup

% mdtable:T10:end
% teaching: NEW-U14-VOC-GODWEBB
% teaching: NEW-U14-NOM-GODWEBB
\reserve{11\baselineskip}
% mdparagraph:P029:begin
\textbf{gōdwębb} — neuter; purple cloth. Stem \textbf{gōdwębb-}. Plural reference; singular productive.
% mdparagraph:P029:end

% table: NEW-U14-NOM-GODWEBB
% mdtable:T11:begin
\par\begingroup
\fontsize{10.5}{13.2}\selectfont
\setlength{\LTleft}{0pt}\setlength{\LTright}{\fill}
\noindent
\begin{tabular}{@{}>{\raggedright\arraybackslash}p{39.46mm}>{\raggedright\arraybackslash}p{39.46mm}>{\raggedright\arraybackslash}p{39.46mm}@{}}
\toprule
\textbf{Case} & \textbf{Singular} & \textbf{Plural} \\
\midrule
Nominative & gōdwębb & gōdwębb \\
Genitive & gōdwębbes & gōdwębba \\
Dative & gōdwębbe & gōdwębbum \\
Accusative & gōdwębb & gōdwębb \\
\bottomrule
\end{tabular}
\par\endgroup

% mdtable:T11:end
% teaching: NEW-U14-VOC-WIC
% teaching: NEW-U14-NOM-WIC
\reserve{11\baselineskip}
% mdparagraph:P030:begin
\textbf{wīc} — neuter; dwelling; abode. Stem \textbf{wīc-}.
% mdparagraph:P030:end

% table: NEW-U14-NOM-WIC
% mdtable:T12:begin
\par\begingroup
\fontsize{10.5}{13.2}\selectfont
\setlength{\LTleft}{0pt}\setlength{\LTright}{\fill}
\noindent
\begin{tabular}{@{}>{\raggedright\arraybackslash}p{39.46mm}>{\raggedright\arraybackslash}p{39.46mm}>{\raggedright\arraybackslash}p{39.46mm}@{}}
\toprule
\textbf{Case} & \textbf{Singular} & \textbf{Plural} \\
\midrule
Nominative & wīc & wīc \\
Genitive & wīces & wīca \\
Dative & wīce & wīcum \\
Accusative & wīc & wīc \\
\bottomrule
\end{tabular}
\par\endgroup

% mdtable:T12:end
% teaching: NEW-U14-VOC-GOD-GOODS
% teaching: NEW-U14-NOM-GOD-GOODS
\reserve{12\baselineskip}
% mdparagraph:P031:begin
\textbf{gōd} — neuter; goods; good things; property (plural). Stem \textbf{gōd-}. Singular reference; plural productive in the goods/\allowbreak{}property sense.
% mdparagraph:P031:end

% table: NEW-U14-NOM-GOD-GOODS
% mdtable:T13:begin
\par\begingroup
\fontsize{10.5}{13.2}\selectfont
\setlength{\LTleft}{0pt}\setlength{\LTright}{\fill}
\noindent
\begin{tabular}{@{}>{\raggedright\arraybackslash}p{39.46mm}>{\raggedright\arraybackslash}p{39.46mm}>{\raggedright\arraybackslash}p{39.46mm}@{}}
\toprule
\textbf{Case} & \textbf{Singular} & \textbf{Plural} \\
\midrule
Nominative & gōd & gōd \\
Genitive & gōdes & gōda \\
Dative & gōde & gōdum \\
Accusative & gōd & gōd \\
\bottomrule
\end{tabular}
\par\endgroup

% mdtable:T13:end
\reserve{7\baselineskip}
% mdparagraph:P032:begin
\subsection*{The reduced stem of wuldor}
% mdparagraph:P032:end

% mdparagraph:P033:begin
Wuldor loses the unstressed o before an ending: wuldr-. Its reference nominative/\allowbreak{}accusative plural wuldru follows the admitted disyllabic-neuter -u pattern. That plural is a rule-based reference form; it does not assign a new count meaning of “glories”.
% mdparagraph:P033:end

% teaching: NEW-U14-VOC-WULDOR
% teaching: NEW-U14-NOM-WULDOR
\reserve{11\baselineskip}
% mdparagraph:P034:begin
\textbf{wuldor} — neuter; glory. Stem \textbf{wuldr-}. Plural reference; singular productive.
% mdparagraph:P034:end

% table: NEW-U14-NOM-WULDOR
% mdtable:T14:begin
\par\begingroup
\fontsize{10.5}{13.2}\selectfont
\setlength{\LTleft}{0pt}\setlength{\LTright}{\fill}
\noindent
\begin{tabular}{@{}>{\raggedright\arraybackslash}p{39.46mm}>{\raggedright\arraybackslash}p{39.46mm}>{\raggedright\arraybackslash}p{39.46mm}@{}}
\toprule
\textbf{Case} & \textbf{Singular} & \textbf{Plural} \\
\midrule
Nominative & wuldor & wuldru \\
Genitive & wuldres & wuldra \\
Dative & wuldre & wuldrum \\
Accusative & wuldor & wuldru \\
\bottomrule
\end{tabular}
\par\endgroup

% mdtable:T14:end
\reserve{7\baselineskip}
% mdparagraph:P035:begin
\subsection*{Strong feminine nouns}
% mdparagraph:P035:end

% mdparagraph:P036:begin
The nominative singular has no ending; genitive, dative and accusative singular take -e. Nominative, genitive and accusative plural take -a, and dative plural -um. In sōðfæstnes the inflectional stem has ss: sōðfæstness-. The tables select genitive plural -a.
% mdparagraph:P036:end

% teaching: NEW-U14-VOC-SOTHFAESTNES
% teaching: NEW-U14-NOM-SOTHFAESTNES
\reserve{11\baselineskip}
% mdparagraph:P037:begin
\textbf{sōðfæstnes} — feminine; truth. Stem \textbf{sōðfæstness-}. Plural reference; singular productive.
% mdparagraph:P037:end

% table: NEW-U14-NOM-SOTHFAESTNES
% mdtable:T15:begin
\par\begingroup
\fontsize{10.5}{13.2}\selectfont
\setlength{\LTleft}{0pt}\setlength{\LTright}{\fill}
\noindent
\begin{tabular}{@{}>{\raggedright\arraybackslash}p{33.82mm}>{\raggedright\arraybackslash}p{40.59mm}>{\raggedright\arraybackslash}p{43.97mm}@{}}
\toprule
\textbf{Case} & \textbf{Singular} & \textbf{Plural} \\
\midrule
Nominative & sōðfæstnes & sōðfæstnessa \\
Genitive & sōðfæstnesse & sōðfæstnessa \\
Dative & sōðfæstnesse & sōðfæstnessum \\
Accusative & sōðfæstnesse & sōðfæstnessa \\
\bottomrule
\end{tabular}
\par\endgroup

% mdtable:T15:end
% teaching: NEW-U14-VOC-SIBB
% teaching: NEW-U14-NOM-SIBB
\reserve{11\baselineskip}
% mdparagraph:P038:begin
\textbf{sibb} — feminine; peace. Stem \textbf{sibb-}. Plural reference; singular productive.
% mdparagraph:P038:end

% table: NEW-U14-NOM-SIBB
% mdtable:T16:begin
\par\begingroup
\fontsize{10.5}{13.2}\selectfont
\setlength{\LTleft}{0pt}\setlength{\LTright}{\fill}
\noindent
\begin{tabular}{@{}>{\raggedright\arraybackslash}p{39.46mm}>{\raggedright\arraybackslash}p{39.46mm}>{\raggedright\arraybackslash}p{39.46mm}@{}}
\toprule
\textbf{Case} & \textbf{Singular} & \textbf{Plural} \\
\midrule
Nominative & sibb & sibba \\
Genitive & sibbe & sibba \\
Dative & sibbe & sibbum \\
Accusative & sibbe & sibba \\
\bottomrule
\end{tabular}
\par\endgroup

% mdtable:T16:end
% teaching: NEW-U14-VOC-WUND
% teaching: NEW-U14-NOM-WUND
\reserve{11\baselineskip}
% mdparagraph:P039:begin
\textbf{wund} — feminine; wound; sore. Stem \textbf{wund-}.
% mdparagraph:P039:end

% table: NEW-U14-NOM-WUND
% mdtable:T17:begin
\par\begingroup
\fontsize{10.5}{13.2}\selectfont
\setlength{\LTleft}{0pt}\setlength{\LTright}{\fill}
\noindent
\begin{tabular}{@{}>{\raggedright\arraybackslash}p{39.46mm}>{\raggedright\arraybackslash}p{39.46mm}>{\raggedright\arraybackslash}p{39.46mm}@{}}
\toprule
\textbf{Case} & \textbf{Singular} & \textbf{Plural} \\
\midrule
Nominative & wund & wunda \\
Genitive & wunde & wunda \\
Dative & wunde & wundum \\
Accusative & wunde & wunda \\
\bottomrule
\end{tabular}
\par\endgroup

% mdtable:T17:end
% mdparagraph:P040:begin
This feminine noun means “wound; sore”; it is distinct from the adjective “wounded”. See correction C04 in the editorial notes.
% mdparagraph:P040:end

\reserve{7\baselineskip}
% mdparagraph:P041:begin
\subsection*{Weak nouns}
% mdparagraph:P041:end

% mdparagraph:P042:begin
The weak masculine citation ends in -a; its other singular cases and its nominative/\allowbreak{}accusative plural end in -an. The weak neuter ēage instead keeps -e in both nominative and accusative singular. Its other -an forms coincide with those of the masculine pattern. Thus -an alone cannot decide case or number. All four nouns use genitive plural -ena and dative plural -um; ēage also permits genitive plural ēagna beside selected ēagena.
% mdparagraph:P042:end

% teaching: NEW-U14-VOC-WAEDLA
% teaching: NEW-U14-NOM-WAEDLA
\reserve{11\baselineskip}
% mdparagraph:P043:begin
\textbf{wǣdla} — masculine; poor or destitute man. Stem \textbf{wǣdl-}.
% mdparagraph:P043:end

% table: NEW-U14-NOM-WAEDLA
% mdtable:T18:begin
\par\begingroup
\fontsize{10.5}{13.2}\selectfont
\setlength{\LTleft}{0pt}\setlength{\LTright}{\fill}
\noindent
\begin{tabular}{@{}>{\raggedright\arraybackslash}p{39.46mm}>{\raggedright\arraybackslash}p{39.46mm}>{\raggedright\arraybackslash}p{39.46mm}@{}}
\toprule
\textbf{Case} & \textbf{Singular} & \textbf{Plural} \\
\midrule
Nominative & wǣdla & wǣdlan \\
Genitive & wǣdlan & wǣdlena \\
Dative & wǣdlan & wǣdlum \\
Accusative & wǣdlan & wǣdlan \\
\bottomrule
\end{tabular}
\par\endgroup

% mdtable:T18:end
% teaching: NEW-U14-VOC-WITGA
% teaching: NEW-U14-NOM-WITGA
\reserve{11\baselineskip}
% mdparagraph:P044:begin
\textbf{wītga} — masculine; prophet. Stem \textbf{wītg-}.
% mdparagraph:P044:end

% table: NEW-U14-NOM-WITGA
% mdtable:T19:begin
\par\begingroup
\fontsize{10.5}{13.2}\selectfont
\setlength{\LTleft}{0pt}\setlength{\LTright}{\fill}
\noindent
\begin{tabular}{@{}>{\raggedright\arraybackslash}p{39.46mm}>{\raggedright\arraybackslash}p{39.46mm}>{\raggedright\arraybackslash}p{39.46mm}@{}}
\toprule
\textbf{Case} & \textbf{Singular} & \textbf{Plural} \\
\midrule
Nominative & wītga & wītgan \\
Genitive & wītgan & wītgena \\
Dative & wītgan & wītgum \\
Accusative & wītgan & wītgan \\
\bottomrule
\end{tabular}
\par\endgroup

% mdtable:T19:end
% teaching: NEW-U14-VOC-THEARFA
% teaching: NEW-U14-NOM-THEARFA
\reserve{11\baselineskip}
% mdparagraph:P045:begin
\textbf{ðearfa} — masculine; poor or needy person. Stem \textbf{ðearf-}.
% mdparagraph:P045:end

% table: NEW-U14-NOM-THEARFA
% mdtable:T20:begin
\par\begingroup
\fontsize{10.5}{13.2}\selectfont
\setlength{\LTleft}{0pt}\setlength{\LTright}{\fill}
\noindent
\begin{tabular}{@{}>{\raggedright\arraybackslash}p{39.46mm}>{\raggedright\arraybackslash}p{39.46mm}>{\raggedright\arraybackslash}p{39.46mm}@{}}
\toprule
\textbf{Case} & \textbf{Singular} & \textbf{Plural} \\
\midrule
Nominative & ðearfa & ðearfan \\
Genitive & ðearfan & ðearfena \\
Dative & ðearfan & ðearfum \\
Accusative & ðearfan & ðearfan \\
\bottomrule
\end{tabular}
\par\endgroup

% mdtable:T20:end
% teaching: NEW-U14-VOC-EAGE
% teaching: NEW-U14-NOM-EAGE
\reserve{11\baselineskip}
% mdparagraph:P046:begin
\textbf{ēage} — neuter; eye. Stem \textbf{ēag-}.
% mdparagraph:P046:end

% table: NEW-U14-NOM-EAGE
% mdtable:T21:begin
\par\begingroup
\fontsize{10.5}{13.2}\selectfont
\setlength{\LTleft}{0pt}\setlength{\LTright}{\fill}
\noindent
\begin{tabular}{@{}>{\raggedright\arraybackslash}p{39.46mm}>{\raggedright\arraybackslash}p{39.46mm}>{\raggedright\arraybackslash}p{39.46mm}@{}}
\toprule
\textbf{Case} & \textbf{Singular} & \textbf{Plural} \\
\midrule
Nominative & ēage & ēagan \\
Genitive & ēagan & ēagena \\
Dative & ēagan & ēagum \\
Accusative & ēage & ēagan \\
\bottomrule
\end{tabular}
\par\endgroup

% mdtable:T21:end
\reserve{7\baselineskip}
% mdparagraph:P047:begin
\section*{3. Indicative systems}
% mdparagraph:P047:end

% mdparagraph:P048:begin
The nine systems below give all present and preterite indicative persons. A plural form is repeated for each person to make the complete system explicit. The ordinary meanings and government follow in §5. The four principal parts identify the stems; their first and fourth columns here are diagnostic citations. They do not add productive infinitive or participial uses to these nine systems.\textsuperscript{5}
% mdparagraph:P048:end

\reserve{21\baselineskip}
% mdparagraph:P049:begin
The regular selected present endings are singular -e, -st, -ð and plural -að, with the stem changes and contractions stated beside each table. Weak II has its separate present pattern below. Strong preterites have no ending in first/\allowbreak{}third singular, -e with the plural stem in second singular, and -on in the plural. Weak preterites use the dental stem with singular -e, -est, -e and plural -on.
% mdparagraph:P049:end

% mdtable:T22:begin
\par\begingroup
\fontsize{10.5}{13.2}\selectfont
\setlength{\LTleft}{0pt}\setlength{\LTright}{\fill}
\noindent
\begin{tabular}{@{}>{\raggedright\arraybackslash}p{27.02mm}>{\raggedright\arraybackslash}p{28.55mm}>{\raggedright\arraybackslash}p{23.36mm}>{\raggedright\arraybackslash}p{36.63mm}@{}}
\toprule
\textbf{Infinitive citation} & \textbf{Preterite 1/\allowbreak{}3 singular} & \textbf{Preterite plural} & \textbf{Past-participle citation} \\
\midrule
gewęndan & gewęnde & gewęndon & gewęnd \\
gelīefan & gelīefde & gelīefdon & gelīefed \\
forgiefan & forgeaf & forgēafon & forgiefen \\
āhrēosan & āhrēas & āhruron & āhroren \\
gebiddan & gebæd & gebǣdon & gebeden \\
ætīewan & ætīewde & ætīewdon & ætīewed \\
liccian & liccode & liccodon & liccod \\
āhębban & āhōf & āhōfon & āhafen \\
onfōn & onfēng & onfēngon & onfangen \\
\bottomrule
\end{tabular}
\par\endgroup

% mdtable:T22:end
% teaching: NEW-U14-VOC-GEWENDAN
% teaching: NEW-U14-PAR-GEWENDAN
% teaching: NEW-U14-PAR-GEWENDAN-IND
% teaching: NEW-U14-GRAM-U14-GEWENDAN-INDICATIVE
\par\begin{minipage}{\linewidth}\raggedright
% mdparagraph:P050:begin
\subsection*{gewęndan — return; turn}
% mdparagraph:P050:end

% mdparagraph:P051:begin
Weak I. Use present stem gewęnd-, with d devoiced before the second/\allowbreak{}third singular endings: gewęntst, gewęnt. Preterite gewęnde simplifies the doubled dental; its stem is gewęnd- plus the selected past endings. The lexical ge- is retained.
% mdparagraph:P051:end

% table: NEW-U14-PAR-GEWENDAN-IND
% mdtable:T23:begin
\par\begingroup
\fontsize{10.5}{13.2}\selectfont
\setlength{\LTleft}{0pt}\setlength{\LTright}{\fill}
\noindent
\begin{tabular}{@{}>{\raggedright\arraybackslash}p{46.63mm}>{\raggedright\arraybackslash}p{35.87mm}>{\raggedright\arraybackslash}p{35.87mm}@{}}
\toprule
\textbf{Person/\allowbreak{}number} & \textbf{Present} & \textbf{Preterite} \\
\midrule
1 singular & gewęnde & gewęnde \\
2 singular & gewęntst & gewęndest \\
3 singular & gewęnt & gewęnde \\
1 plural & gewęndað & gewęndon \\
2 plural & gewęndað & gewęndon \\
3 plural & gewęndað & gewęndon \\
\bottomrule
\end{tabular}
\par\endgroup

% mdtable:T23:end
\end{minipage}\par\addvspace{4pt}
% teaching: NEW-U14-VOC-GELIEFAN
% teaching: NEW-U14-PAR-GELIEFAN
% teaching: NEW-U14-PAR-GELIEFAN-IND
% teaching: NEW-U14-GRAM-U14-GELIEFAN-INDICATIVE
\par\begin{minipage}{\linewidth}\raggedright
% mdparagraph:P052:begin
\subsection*{gelīefan — believe}
% mdparagraph:P052:end

% mdparagraph:P053:begin
Weak I. Retain long īe throughout. Present endings are -e, -st, -ð and plural -að; the dental past stem is gelīefd-, with -e, -est, -e and plural -on. The infinitive and past-participle citations differ from the finite forms.
% mdparagraph:P053:end

% table: NEW-U14-PAR-GELIEFAN-IND
% mdtable:T24:begin
\par\begingroup
\fontsize{10.5}{13.2}\selectfont
\setlength{\LTleft}{0pt}\setlength{\LTright}{\fill}
\noindent
\begin{tabular}{@{}>{\raggedright\arraybackslash}p{46.63mm}>{\raggedright\arraybackslash}p{35.87mm}>{\raggedright\arraybackslash}p{35.87mm}@{}}
\toprule
\textbf{Person/\allowbreak{}number} & \textbf{Present} & \textbf{Preterite} \\
\midrule
1 singular & gelīefe & gelīefde \\
2 singular & gelīefst & gelīefdest \\
3 singular & gelīefð & gelīefde \\
1 plural & gelīefað & gelīefdon \\
2 plural & gelīefað & gelīefdon \\
3 plural & gelīefað & gelīefdon \\
\bottomrule
\end{tabular}
\par\endgroup

% mdtable:T24:end
\end{minipage}\par\addvspace{4pt}
% teaching: NEW-U14-VOC-FORGIEFAN
% teaching: NEW-U14-PAR-FORGIEFAN
% teaching: NEW-U14-PAR-FORGIEFAN-IND
% teaching: NEW-U14-GRAM-U14-FORGIEFAN-INDICATIVE
\par\begin{minipage}{\linewidth}\raggedright
% mdparagraph:P054:begin
\subsection*{forgiefan — give; grant}
% mdparagraph:P054:end

% mdparagraph:P055:begin
Strong V. Principal stems are forgief-, forgeaf-, forgēaf-, forgief-. The selected present forms retain ie throughout. Past second singular uses the plural stem and -e; first and third singular share forgeaf.
% mdparagraph:P055:end

% table: NEW-U14-PAR-FORGIEFAN-IND
% mdtable:T25:begin
\par\begingroup
\fontsize{10.5}{13.2}\selectfont
\setlength{\LTleft}{0pt}\setlength{\LTright}{\fill}
\noindent
\begin{tabular}{@{}>{\raggedright\arraybackslash}p{46.63mm}>{\raggedright\arraybackslash}p{35.87mm}>{\raggedright\arraybackslash}p{35.87mm}@{}}
\toprule
\textbf{Person/\allowbreak{}number} & \textbf{Present} & \textbf{Preterite} \\
\midrule
1 singular & forgiefe & forgeaf \\
2 singular & forgiefst & forgēafe \\
3 singular & forgiefð & forgeaf \\
1 plural & forgiefað & forgēafon \\
2 plural & forgiefað & forgēafon \\
3 plural & forgiefað & forgēafon \\
\bottomrule
\end{tabular}
\par\endgroup

% mdtable:T25:end
\end{minipage}\par\addvspace{4pt}
% teaching: NEW-U14-VOC-AHREOSAN
% teaching: NEW-U14-PAR-AHREOSAN
% teaching: NEW-U14-PAR-AHREOSAN-IND
% teaching: NEW-U14-GRAM-U14-AHREOSAN-INDICATIVE
\par\begin{minipage}{\linewidth}\raggedright
% mdparagraph:P056:begin
\subsection*{āhrēosan — fall; fall upon}
% mdparagraph:P056:end

% mdparagraph:P057:begin
Strong II. The principal stems are āhrēos-, āhrēas-, āhrur-, āhror-. Present second/\allowbreak{}third singular have long īe and the selected contracted ending -st. Past second singular and all plural persons use āhrur-, including the s/\allowbreak{}r alternation.
% mdparagraph:P057:end

% table: NEW-U14-PAR-AHREOSAN-IND
% mdtable:T26:begin
\par\begingroup
\fontsize{10.5}{13.2}\selectfont
\setlength{\LTleft}{0pt}\setlength{\LTright}{\fill}
\noindent
\begin{tabular}{@{}>{\raggedright\arraybackslash}p{46.63mm}>{\raggedright\arraybackslash}p{35.87mm}>{\raggedright\arraybackslash}p{35.87mm}@{}}
\toprule
\textbf{Person/\allowbreak{}number} & \textbf{Present} & \textbf{Preterite} \\
\midrule
1 singular & āhrēose & āhrēas \\
2 singular & āhrīest & āhrure \\
3 singular & āhrīest & āhrēas \\
1 plural & āhrēosað & āhruron \\
2 plural & āhrēosað & āhruron \\
3 plural & āhrēosað & āhruron \\
\bottomrule
\end{tabular}
\par\endgroup

% mdtable:T26:end
\end{minipage}\par\addvspace{4pt}
% teaching: NEW-U14-VOC-GEBIDDAN
% teaching: NEW-U14-PAR-GEBIDDAN
% teaching: NEW-U14-PAR-GEBIDDAN-IND
% teaching: NEW-U14-GRAM-U14-GEBIDDAN-INDICATIVE
\par\begin{minipage}{\linewidth}\raggedright
% mdparagraph:P058:begin
\subsection*{gebiddan — pray}
% mdparagraph:P058:end

% mdparagraph:P059:begin
Strong V. Principal stems are gebidd-, gebæd-, gebǣd-, gebed-. The present second/\allowbreak{}third singular simplify the consonants to gebitst and gebitt. Past second singular uses the long-vowel plural stem gebǣd- with -e.
% mdparagraph:P059:end

% table: NEW-U14-PAR-GEBIDDAN-IND
% mdtable:T27:begin
\par\begingroup
\fontsize{10.5}{13.2}\selectfont
\setlength{\LTleft}{0pt}\setlength{\LTright}{\fill}
\noindent
\begin{tabular}{@{}>{\raggedright\arraybackslash}p{46.63mm}>{\raggedright\arraybackslash}p{35.87mm}>{\raggedright\arraybackslash}p{35.87mm}@{}}
\toprule
\textbf{Person/\allowbreak{}number} & \textbf{Present} & \textbf{Preterite} \\
\midrule
1 singular & gebidde & gebæd \\
2 singular & gebitst & gebǣde \\
3 singular & gebitt & gebæd \\
1 plural & gebiddað & gebǣdon \\
2 plural & gebiddað & gebǣdon \\
3 plural & gebiddað & gebǣdon \\
\bottomrule
\end{tabular}
\par\endgroup

% mdtable:T27:end
\end{minipage}\par\addvspace{4pt}
% teaching: NEW-U14-VOC-AETIEWAN
% teaching: NEW-U14-PAR-AETIEWAN
% teaching: NEW-U14-PAR-AETIEWAN-IND
% teaching: NEW-U14-GRAM-U14-AETIEWAN-INDICATIVE
\par\begin{minipage}{\linewidth}\raggedright
% mdparagraph:P060:begin
\subsection*{ætīewan — show; display; reveal}
% mdparagraph:P060:end

% mdparagraph:P061:begin
Weak I. Retain w: present ætīew-, dental past ætīewd-. This transitive verb “show” is distinct from earlier weak-II ætēowian, “appear”.
% mdparagraph:P061:end

% table: NEW-U14-PAR-AETIEWAN-IND
% mdtable:T28:begin
\par\begingroup
\fontsize{10.5}{13.2}\selectfont
\setlength{\LTleft}{0pt}\setlength{\LTright}{\fill}
\noindent
\begin{tabular}{@{}>{\raggedright\arraybackslash}p{46.63mm}>{\raggedright\arraybackslash}p{35.87mm}>{\raggedright\arraybackslash}p{35.87mm}@{}}
\toprule
\textbf{Person/\allowbreak{}number} & \textbf{Present} & \textbf{Preterite} \\
\midrule
1 singular & ætīewe & ætīewde \\
2 singular & ætīewst & ætīewdest \\
3 singular & ætīewð & ætīewde \\
1 plural & ætīewað & ætīewdon \\
2 plural & ætīewað & ætīewdon \\
3 plural & ætīewað & ætīewdon \\
\bottomrule
\end{tabular}
\par\endgroup

% mdtable:T28:end
\end{minipage}\par\addvspace{4pt}
% teaching: NEW-U14-VOC-LICCIAN
% teaching: NEW-U14-PAR-LICCIAN
% teaching: NEW-U14-PAR-LICCIAN-IND
% teaching: NEW-U14-GRAM-U14-LICCIAN-INDICATIVE
\par\begin{minipage}{\linewidth}\raggedright
% mdparagraph:P062:begin
\subsection*{liccian — lick}
% mdparagraph:P062:end

% mdparagraph:P063:begin
Weak II. Present first singular has -ie, second -ast, third -að and plural -iað. The past stem liccod- takes -e, -est, -e and plural -on. Both c letters remain; the selected past uses -ode, not an unassigned alternative in -ade.
% mdparagraph:P063:end

% table: NEW-U14-PAR-LICCIAN-IND
% mdtable:T29:begin
\par\begingroup
\fontsize{10.5}{13.2}\selectfont
\setlength{\LTleft}{0pt}\setlength{\LTright}{\fill}
\noindent
\begin{tabular}{@{}>{\raggedright\arraybackslash}p{46.63mm}>{\raggedright\arraybackslash}p{35.87mm}>{\raggedright\arraybackslash}p{35.87mm}@{}}
\toprule
\textbf{Person/\allowbreak{}number} & \textbf{Present} & \textbf{Preterite} \\
\midrule
1 singular & liccie & liccode \\
2 singular & liccast & liccodest \\
3 singular & liccað & liccode \\
1 plural & licciað & liccodon \\
2 plural & licciað & liccodon \\
3 plural & licciað & liccodon \\
\bottomrule
\end{tabular}
\par\endgroup

% mdtable:T29:end
\end{minipage}\par\addvspace{4pt}
% teaching: NEW-U14-VOC-AHEBBAN
% teaching: NEW-U14-PAR-AHEBBAN
% teaching: NEW-U14-PAR-AHEBBAN-IND
% teaching: NEW-U14-GRAM-U14-AHEBBAN-INDICATIVE
\par\begin{minipage}{\linewidth}\raggedright
% mdparagraph:P064:begin
\subsection*{āhębban — raise; lift}
% mdparagraph:P064:end

% mdparagraph:P065:begin
Strong VI. Present āhębb- appears with bb in first singular and plural, but with f before second/\allowbreak{}third singular endings. The past stem is āhōf-: zero ending in first/\allowbreak{}third singular, -e in second singular and -on in the plural. Retain the prefix ā-.
% mdparagraph:P065:end

% table: NEW-U14-PAR-AHEBBAN-IND
% mdtable:T30:begin
\par\begingroup
\fontsize{10.5}{13.2}\selectfont
\setlength{\LTleft}{0pt}\setlength{\LTright}{\fill}
\noindent
\begin{tabular}{@{}>{\raggedright\arraybackslash}p{46.63mm}>{\raggedright\arraybackslash}p{35.87mm}>{\raggedright\arraybackslash}p{35.87mm}@{}}
\toprule
\textbf{Person/\allowbreak{}number} & \textbf{Present} & \textbf{Preterite} \\
\midrule
1 singular & āhębbe & āhōf \\
2 singular & āhęfst & āhōfe \\
3 singular & āhęfð & āhōf \\
1 plural & āhębbað & āhōfon \\
2 plural & āhębbað & āhōfon \\
3 plural & āhębbað & āhōfon \\
\bottomrule
\end{tabular}
\par\endgroup

% mdtable:T30:end
\end{minipage}\par\addvspace{4pt}
% teaching: NEW-U14-VOC-ONFON
% teaching: NEW-U14-PAR-ONFON
% teaching: NEW-U14-PAR-ONFON-IND
% teaching: NEW-U14-GRAM-U14-ONFON-INDICATIVE
\par\begin{minipage}{\linewidth}\raggedright
% mdparagraph:P066:begin
\subsection*{onfōn — receive; accept}
% mdparagraph:P066:end

% mdparagraph:P067:begin
Contracted strong VII. Compare present onfō, onfēhst, onfēhð and plural onfōð. The past stem onfēng- takes no ending in first/\allowbreak{}third singular, -e in second singular and -on in all plural persons. Its government must be learned separately.
% mdparagraph:P067:end

% table: NEW-U14-PAR-ONFON-IND
% mdtable:T31:begin
\par\begingroup
\fontsize{10.5}{13.2}\selectfont
\setlength{\LTleft}{0pt}\setlength{\LTright}{\fill}
\noindent
\begin{tabular}{@{}>{\raggedright\arraybackslash}p{46.63mm}>{\raggedright\arraybackslash}p{35.87mm}>{\raggedright\arraybackslash}p{35.87mm}@{}}
\toprule
\textbf{Person/\allowbreak{}number} & \textbf{Present} & \textbf{Preterite} \\
\midrule
1 singular & onfō & onfēng \\
2 singular & onfēhst & onfēnge \\
3 singular & onfēhð & onfēng \\
1 plural & onfōð & onfēngon \\
2 plural & onfōð & onfēngon \\
3 plural & onfōð & onfēngon \\
\bottomrule
\end{tabular}
\par\endgroup

% mdtable:T31:end
\end{minipage}\par\addvspace{4pt}
\reserve{7\baselineskip}
% mdparagraph:P068:begin
\section*{4. Subjunctives, participles and commands}
% mdparagraph:P068:end

\reserve{7\baselineskip}
% mdparagraph:P069:begin
\subsection*{Five complete subjunctive systems}
% mdparagraph:P069:end

% mdparagraph:P070:begin
In each tense the subjunctive has one singular form for all three persons, ending in -e, and one plural form, ending in -en. Strong preterite subjunctives use the preterite plural stem; weak ones retain the dental past stem. Syncretism does not remove a distinction of tense or mood: sęnde and gelīefde require their clause context. The already taught purpose and request clauses select subjunctives; a factual ðæt-clause need not do so. The indefinite concessive expression in Cook 27.10 receives separate local help.
% mdparagraph:P070:end

% teaching: NEW-U14-VOC-FORWEORTHAN
% teaching: NEW-U14-PAR-FORWEORTHAN
% teaching: NEW-U14-PAR-FORWEORTHAN-SUBJ
% teaching: NEW-U14-GRAM-U14-FORWEORTHAN-SUBJUNCTIVE
\reserve{13\baselineskip}
% mdparagraph:P071:begin
\textbf{forweorðan.} “Perish”. Present stem forweorð-; past forwurd-. Keep the ð/\allowbreak{}d alternation and the short past u.
% mdparagraph:P071:end

% table: NEW-U14-PAR-FORWEORTHAN-SUBJ
% mdtable:T32:begin
\par\begingroup
\fontsize{10.5}{13.2}\selectfont
\setlength{\LTleft}{0pt}\setlength{\LTright}{\fill}
\noindent
\begin{tabular}{@{}>{\raggedright\arraybackslash}p{46.63mm}>{\raggedright\arraybackslash}p{35.87mm}>{\raggedright\arraybackslash}p{35.87mm}@{}}
\toprule
\textbf{Person/\allowbreak{}number} & \textbf{Present} & \textbf{Preterite} \\
\midrule
1 singular & forweorðe & forwurde \\
2 singular & forweorðe & forwurde \\
3 singular & forweorðe & forwurde \\
1 plural & forweorðen & forwurden \\
2 plural & forweorðen & forwurden \\
3 plural & forweorðen & forwurden \\
\bottomrule
\end{tabular}
\par\endgroup

% mdtable:T32:end
% teaching: NEW-U14-PAR-GEBIDDAN-SUBJ
% teaching: NEW-U14-GRAM-U14-GEBIDDAN-SUBJUNCTIVE
\reserve{14\baselineskip}
% mdparagraph:P072:begin
\textbf{gebiddan.} “Pray”. Present stem gebidd-; past gebǣd-. The accusative reflexive is retained in the purpose use as well as in factual indicative clauses.
% mdparagraph:P072:end

% table: NEW-U14-PAR-GEBIDDAN-SUBJ
% mdtable:T33:begin
\par\begingroup
\fontsize{10.5}{13.2}\selectfont
\setlength{\LTleft}{0pt}\setlength{\LTright}{\fill}
\noindent
\begin{tabular}{@{}>{\raggedright\arraybackslash}p{46.63mm}>{\raggedright\arraybackslash}p{35.87mm}>{\raggedright\arraybackslash}p{35.87mm}@{}}
\toprule
\textbf{Person/\allowbreak{}number} & \textbf{Present} & \textbf{Preterite} \\
\midrule
1 singular & gebidde & gebǣde \\
2 singular & gebidde & gebǣde \\
3 singular & gebidde & gebǣde \\
1 plural & gebidden & gebǣden \\
2 plural & gebidden & gebǣden \\
3 plural & gebidden & gebǣden \\
\bottomrule
\end{tabular}
\par\endgroup

% mdtable:T33:end
% teaching: NEW-U14-VOC-SENDAN
% teaching: NEW-U14-PAR-SENDAN
% teaching: NEW-U14-PAR-SENDAN-SUBJ
% teaching: NEW-U14-GRAM-U14-SENDAN-SUBJUNCTIVE
\reserve{15\baselineskip}
% mdparagraph:P073:begin
\textbf{sęndan.} “Send”. The dental past simplifies -ndd- to -nd-, so present and past forms coincide in the selected subjunctive. The earlier recipient frame remains accusative person or thing sent, dative recipient.
% mdparagraph:P073:end

% table: NEW-U14-PAR-SENDAN-SUBJ
% mdtable:T34:begin
\par\begingroup
\fontsize{10.5}{13.2}\selectfont
\setlength{\LTleft}{0pt}\setlength{\LTright}{\fill}
\noindent
\begin{tabular}{@{}>{\raggedright\arraybackslash}p{46.63mm}>{\raggedright\arraybackslash}p{35.87mm}>{\raggedright\arraybackslash}p{35.87mm}@{}}
\toprule
\textbf{Person/\allowbreak{}number} & \textbf{Present} & \textbf{Preterite} \\
\midrule
1 singular & sęnde & sęnde \\
2 singular & sęnde & sęnde \\
3 singular & sęnde & sęnde \\
1 plural & sęnden & sęnden \\
2 plural & sęnden & sęnden \\
3 plural & sęnden & sęnden \\
\bottomrule
\end{tabular}
\par\endgroup

% mdtable:T34:end
% teaching: NEW-U14-PAR-GELIEFAN-SUBJ
% teaching: NEW-U14-GRAM-U14-GELIEFAN-SUBJUNCTIVE
\reserve{13\baselineskip}
% mdparagraph:P074:begin
\textbf{gelīefan.} “Believe”. Present gelīef-; past gelīefd-. Keep the dental and long īe.
% mdparagraph:P074:end

% table: NEW-U14-PAR-GELIEFAN-SUBJ
% mdtable:T35:begin
\par\begingroup
\fontsize{10.5}{13.2}\selectfont
\setlength{\LTleft}{0pt}\setlength{\LTright}{\fill}
\noindent
\begin{tabular}{@{}>{\raggedright\arraybackslash}p{46.63mm}>{\raggedright\arraybackslash}p{35.87mm}>{\raggedright\arraybackslash}p{35.87mm}@{}}
\toprule
\textbf{Person/\allowbreak{}number} & \textbf{Present} & \textbf{Preterite} \\
\midrule
1 singular & gelīefe & gelīefde \\
2 singular & gelīefe & gelīefde \\
3 singular & gelīefe & gelīefde \\
1 plural & gelīefen & gelīefden \\
2 plural & gelīefen & gelīefden \\
3 plural & gelīefen & gelīefden \\
\bottomrule
\end{tabular}
\par\endgroup

% mdtable:T35:end
% teaching: NEW-U14-VOC-ARISAN
% teaching: NEW-U14-PAR-ARISAN
% teaching: NEW-U14-PAR-ARISAN-SUBJ
% teaching: NEW-U14-GRAM-U14-ARISAN-SUBJUNCTIVE
\reserve{13\baselineskip}
% mdparagraph:P075:begin
\textbf{ārīsan.} “Arise”. Present ārīs- has long ī; past āris- has short i. The prefix ā- stays long in both tenses.
% mdparagraph:P075:end

% table: NEW-U14-PAR-ARISAN-SUBJ
% mdtable:T36:begin
\par\begingroup
\fontsize{10.5}{13.2}\selectfont
\setlength{\LTleft}{0pt}\setlength{\LTright}{\fill}
\noindent
\begin{tabular}{@{}>{\raggedright\arraybackslash}p{46.63mm}>{\raggedright\arraybackslash}p{35.87mm}>{\raggedright\arraybackslash}p{35.87mm}@{}}
\toprule
\textbf{Person/\allowbreak{}number} & \textbf{Present} & \textbf{Preterite} \\
\midrule
1 singular & ārīse & ārise \\
2 singular & ārīse & ārise \\
3 singular & ārīse & ārise \\
1 plural & ārīsen & ārisen \\
2 plural & ārīsen & ārisen \\
3 plural & ārīsen & ārisen \\
\bottomrule
\end{tabular}
\par\endgroup

% mdtable:T36:end
\reserve{7\baselineskip}
% mdparagraph:P076:begin
\subsection*{Four participle-base pairs}
% mdparagraph:P076:end

% mdparagraph:P077:begin
The present base uses the present stem + -ende. A strong past base uses its fourth stem + -en; a weak past base contains the dental formation. These four pairs teach the bases and the selected past-passive use. They do not establish complete attributive participle declensions. The source’s inflected present-participial phrases are supplied separately.
% mdparagraph:P077:end

% teaching: NEW-U14-VOC-FORCEORFAN
% teaching: NEW-U14-PAR-FORCEORFAN
% teaching: NEW-U14-PAR-FORCEORFAN-PART
% teaching: NEW-U14-GRAM-U14-FORCEORFAN-PARTICIPLE-BASES
\reserve{9\baselineskip}
% mdparagraph:P078:begin
\textbf{forceorfan.} “Cut down; cut off”. Present forceorf-; strong past forcorf- + -en.
% mdparagraph:P078:end

% table: NEW-U14-PAR-FORCEORFAN-PART
% mdtable:T37:begin
\par\begingroup
\fontsize{10.5}{13.2}\selectfont
\setlength{\LTleft}{0pt}\setlength{\LTright}{\fill}
\noindent
\begin{tabular}{@{}>{\raggedright\arraybackslash}p{68.97mm}>{\raggedright\arraybackslash}p{52.22mm}@{}}
\toprule
\textbf{Cell} & \textbf{Form} \\
\midrule
Present participle base & forceorfende \\
Past participle base & forcorfen \\
\bottomrule
\end{tabular}
\par\endgroup

% mdtable:T37:end
% teaching: NEW-U14-VOC-WEORPAN
% teaching: NEW-U14-PAR-WEORPAN
% teaching: NEW-U14-PAR-WEORPAN-PART
% teaching: NEW-U14-GRAM-U14-WEORPAN-PARTICIPLE-BASES
\reserve{10\baselineskip}
% mdparagraph:P079:begin
\textbf{weorpan.} “Throw”. Present weorp-; strong past worp- + -en. Selected worpen is valid without ge-; geworpen is an accepted alternative.
% mdparagraph:P079:end

% table: NEW-U14-PAR-WEORPAN-PART
% mdtable:T38:begin
\par\begingroup
\fontsize{10.5}{13.2}\selectfont
\setlength{\LTleft}{0pt}\setlength{\LTright}{\fill}
\noindent
\begin{tabular}{@{}>{\raggedright\arraybackslash}p{74.31mm}>{\raggedright\arraybackslash}p{46.88mm}@{}}
\toprule
\textbf{Cell} & \textbf{Form} \\
\midrule
Present participle base & weorpende \\
Past participle base & worpen \\
\bottomrule
\end{tabular}
\par\endgroup

% mdtable:T38:end
% teaching: NEW-U14-VOC-GEMETAN
% teaching: NEW-U14-PAR-GEMETAN
% teaching: NEW-U14-PAR-GEMETAN-PART
% teaching: NEW-U14-GRAM-U14-GEMETAN-PARTICIPLE-BASES
\reserve{9\baselineskip}
% mdparagraph:P080:begin
\textbf{gemētan.} “Find; meet”. Present gemēt-; the dental past participle contracts to gemētt. Retain the lexical ge- once.
% mdparagraph:P080:end

% table: NEW-U14-PAR-GEMETAN-PART
% mdtable:T39:begin
\par\begingroup
\fontsize{10.5}{13.2}\selectfont
\setlength{\LTleft}{0pt}\setlength{\LTright}{\fill}
\noindent
\begin{tabular}{@{}>{\raggedright\arraybackslash}p{74.31mm}>{\raggedright\arraybackslash}p{46.88mm}@{}}
\toprule
\textbf{Cell} & \textbf{Form} \\
\midrule
Present participle base & gemētende \\
Past participle base & gemētt \\
\bottomrule
\end{tabular}
\par\endgroup

% mdtable:T39:end
% teaching: NEW-U14-VOC-GEGLENGAN
% teaching: NEW-U14-PAR-GEGLENGAN
% teaching: NEW-U14-PAR-GEGLENGAN-PART
% teaching: NEW-U14-GRAM-U14-GEGLENGAN-PARTICIPLE-BASES
\reserve{9\baselineskip}
% mdparagraph:P081:begin
\textbf{geglęngan.} “Adorn”. Present geglęng-; weak past base geglęnged. Do not add a second ge-.
% mdparagraph:P081:end

% table: NEW-U14-PAR-GEGLENGAN-PART
% mdtable:T40:begin
\par\begingroup
\fontsize{10.5}{13.2}\selectfont
\setlength{\LTleft}{0pt}\setlength{\LTright}{\fill}
\noindent
\begin{tabular}{@{}>{\raggedright\arraybackslash}p{71.54mm}>{\raggedright\arraybackslash}p{49.65mm}@{}}
\toprule
\textbf{Cell} & \textbf{Form} \\
\midrule
Present participle base & geglęngende \\
Past participle base & geglęnged \\
\bottomrule
\end{tabular}
\par\endgroup

% mdtable:T40:end
% mdparagraph:P082:begin
With the earlier passive construction, the person or thing affected is the subject and the finite form of bēon or wesan carries tense. Coordinated past participles may share one auxiliary. An instruction “verb last” in Cook 26.7 or 27.2 selects a final finite auxiliary after the participle group and its complements; it does not merely put a nonfinite participle last.
% mdparagraph:P082:end

\reserve{7\baselineskip}
% mdparagraph:P083:begin
\subsection*{Commands}
% mdparagraph:P083:end

% mdparagraph:P084:begin
These weak-II commands have singular -a and plural -iað. The noun or pronoun after the verb takes the case required by its own construction. Preserve ts in gemiltsian and the single l in tilian.
% mdparagraph:P084:end

% teaching: NEW-U14-VOC-GEMILTSIAN
% teaching: NEW-U14-PAR-GEMILTSIAN
% teaching: NEW-U14-PAR-GEMILTSIAN-IMP
% teaching: NEW-U14-GRAM-U14-GEMILTSIAN-IMPERATIVE
\reserve{9\baselineskip}
% mdparagraph:P085:begin
\textbf{gemiltsian} — have compassion; have mercy.
% mdparagraph:P085:end

% table: NEW-U14-PAR-GEMILTSIAN-IMP
% mdtable:T41:begin
\par\begingroup
\fontsize{10.5}{13.2}\selectfont
\setlength{\LTleft}{0pt}\setlength{\LTright}{\fill}
\noindent
\begin{tabular}{@{}>{\raggedright\arraybackslash}p{65.08mm}>{\raggedright\arraybackslash}p{56.11mm}@{}}
\toprule
\textbf{Cell} & \textbf{Form} \\
\midrule
Singular command & gemiltsa \\
Plural command & gemiltsiað \\
\bottomrule
\end{tabular}
\par\endgroup

% mdtable:T41:end
% teaching: NEW-U14-VOC-TILIAN
% teaching: NEW-U14-PAR-TILIAN
% teaching: NEW-U14-PAR-TILIAN-IMP
% teaching: NEW-U14-GRAM-U14-TILIAN-IMPERATIVE
\reserve{9\baselineskip}
% mdparagraph:P086:begin
\textbf{tilian} — obtain; provide.
% mdparagraph:P086:end

% table: NEW-U14-PAR-TILIAN-IMP
% mdtable:T42:begin
\par\begingroup
\fontsize{10.5}{13.2}\selectfont
\setlength{\LTleft}{0pt}\setlength{\LTright}{\fill}
\noindent
\begin{tabular}{@{}>{\raggedright\arraybackslash}p{65.08mm}>{\raggedright\arraybackslash}p{56.11mm}@{}}
\toprule
\textbf{Cell} & \textbf{Form} \\
\midrule
Singular command & tila \\
Plural command & tiliað \\
\bottomrule
\end{tabular}
\par\endgroup

% mdtable:T42:end
\reserve{7\baselineskip}
% mdparagraph:P087:begin
\section*{5. Meanings and ordinary uses}
% mdparagraph:P087:end

% mdparagraph:P088:begin
A verb’s forms and its government are separate facts. The following uses apply with the systems assigned above and earlier. In particular, prefixes do not preserve a base verb’s government automatically.
% mdparagraph:P088:end

% teaching: NEW-U14-VOC-FORLAETAN
% teaching: NEW-U14-VOC-ON
% teaching: NEW-U14-PAR-ON
% teaching: NEW-U14-VOC-MID
% teaching: NEW-U14-PAR-MID
% teaching: NEW-U14-GRAM-U14-GEWENDAN-RETURN
% teaching: NEW-U14-GRAM-U14-AHREOSAN-ON
% teaching: NEW-U14-GRAM-U14-GAN-ON-GOAL
% teaching: NEW-U14-GRAM-U14-BERAN-TRANSPORT
% teaching: NEW-U14-GRAM-U14-ARISAN-SOURCE
% mdparagraph:P089:begin
\textbf{Motion and destination.} Gewęndan, “return”, is used intransitively with tō + dative for the endpoint. Āhrēosan, “fall”, takes on + accusative when falling upon a goal. Gān takes on + accusative for entry or ascent into or onto a place. Beran carries an accusative person or thing to a destination marked tō + dative. Ārīsan is intransitive; of + dative marks the source from which something arises.
% mdparagraph:P089:end

% teaching: NEW-U14-GRAM-U14-GELIEFAN-ON
% mdparagraph:P090:begin
\textbf{Belief.} Gelīefan on + accusative person means “believe in” that person. This is a lexical use of on, not a physical-motion construction. Gelīefan may also stand without an expressed complement, “believe”.
% mdparagraph:P090:end

% teaching: NEW-U14-GRAM-U14-FORLAETAN-TRANSFER
% teaching: NEW-U14-GRAM-U14-FORGIEFAN-GIVING
% teaching: NEW-U14-GRAM-U14-AETIEWAN-SHOW
% teaching: NEW-U14-GRAM-U14-GIEFAN-UNEXPRESSED-THING
% mdparagraph:P091:begin
\textbf{Giving, leaving and showing.} Forgiefan here means “give; grant”, with an accusative thing and a dative recipient. Existing forlǣtan has the same cases in the selected expression “leave peace to someone”; its earlier abandonment meaning alone did not teach this use. Ætīewan shows an accusative thing to a dative person. With giefan the thing given may be left unexpressed when the context supplies it; the recipient remains dative.
% mdparagraph:P091:end

% teaching: NEW-U14-GRAM-U14-GEBIDDAN-REFLEXIVE
% mdparagraph:P092:begin
\textbf{Prayer.} Gebiddan uses an accusative personal pronoun referring to its subject. Keep that reflexive in both the purpose clause and the factual when-clause of Cook 26.8. English normally translates the whole construction simply as “pray”. The purpose clause uses subjunctive; the factual past when-clause uses indicative.
% mdparagraph:P092:end

% teaching: NEW-U14-GRAM-U14-LICCIAN-OBJECT
% teaching: NEW-U14-GRAM-U14-AHEBBAN-OBJECT
% mdparagraph:P093:begin
\textbf{Affected objects.} Liccian, “lick”, takes an accusative object. Āhębban, “raise; lift”, also takes an accusative object; ūp may be added as a separate adverb.
% mdparagraph:P093:end

% teaching: NEW-U14-GRAM-U14-WEORPAN-GOAL-PASSIVE
% teaching: NEW-U14-GRAM-U14-GEGLENGAN-MATERIAL
% mdparagraph:P094:begin
\textbf{Passive complements.} Weorpan throws an accusative thing into or onto an on + accusative goal. Its passive retains that goal. Geglęngan adorns a person; the passive makes that person the subject and may add mid + dative for the material or clothing.
% mdparagraph:P094:end

% teaching: NEW-U14-GRAM-U14-ONFON-DATIVE
% teaching: NEW-U14-GRAM-U14-GEMILTSIAN-DATIVE
% mdparagraph:P095:begin
\textbf{Dative government.} Onfōn, “receive”, takes the dative thing received. This differs from the earlier accusative after fōn, “catch”. Gemiltsian, “have mercy on”, takes the dative person.
% mdparagraph:P095:end

% teaching: NEW-U14-GRAM-U14-TILIAN-GENITIVE-REFLEXIVE
% mdparagraph:P096:begin
\textbf{Providing friends.} In the selected tilian construction the friends acquired are genitive; the personal beneficiary “for yourselves” is a dative reflexive. Both case requirements must be expressed. Singular and plural commands use the same government.
% mdparagraph:P096:end

% mdparagraph:P097:begin
\textbf{Earlier clause uses.} The bare relative particle ðe can introduce a relative clause without a preceding demonstrative inside that clause. The noun’s role in its main clause need not match the relative’s role. The purpose clauses in Cook 26.2 use subjunctives without an added “may” auxiliary. For Cook 27.9, use the earlier hālsian request construction: the addressee is accusative, and a following ðæt-clause expresses the requested action in the subjunctive. Sęndan retains its earlier accusative object and dative recipient where a recipient is expressed.
% mdparagraph:P097:end

\reserve{8\baselineskip}
% example: NEW-U14-EXAMPLE01
\reserve{5\baselineskip}
% mdparagraph:P098:begin
\textbf{Example 1.}
% mdparagraph:P098:end

\reserve{5\baselineskip}
% mdparagraph:P099:begin
\textbf{Ðæt goldhord is on ðǣm hūse.}
% mdparagraph:P099:end

% mdparagraph:P100:begin
The treasure is in the house.
% mdparagraph:P100:end

% mdparagraph:P101:begin
On + dative gives the location of the treasure. The compound is one neuter noun.
% mdparagraph:P101:end

\reserve{8\baselineskip}
% example: NEW-U14-EXAMPLE02
\reserve{5\baselineskip}
% mdparagraph:P102:begin
\textbf{Example 2.}
% mdparagraph:P102:end

\reserve{5\baselineskip}
% mdparagraph:P103:begin
\textbf{Ic forlǣte ēow sibbe.}
% mdparagraph:P103:end

% mdparagraph:P104:begin
I leave you peace.
% mdparagraph:P104:end

% mdparagraph:P105:begin
Sibbe is the thing left; ēow names the plural recipients.
% mdparagraph:P105:end

\reserve{8\baselineskip}
% example: NEW-U14-EXAMPLE03
\reserve{5\baselineskip}
% mdparagraph:P106:begin
\textbf{Example 3.}
% mdparagraph:P106:end

\reserve{5\baselineskip}
% mdparagraph:P107:begin
\textbf{Hē onfēng ðǣm golde.}
% mdparagraph:P107:end

% mdparagraph:P108:begin
He received the gold.
% mdparagraph:P108:end

% mdparagraph:P109:begin
The thing received is dative after onfōn.
% mdparagraph:P109:end

\reserve{8\baselineskip}
% example: NEW-U14-EXAMPLE04
\reserve{5\baselineskip}
% mdparagraph:P110:begin
\textbf{Example 4.}
% mdparagraph:P110:end

\reserve{5\baselineskip}
% mdparagraph:P111:begin
\textbf{Gemiltsa mē.}
% mdparagraph:P111:end

% mdparagraph:P112:begin
Have mercy on me.
% mdparagraph:P112:end

% mdparagraph:P113:begin
The singular command governs dative mē.
% mdparagraph:P113:end

\reserve{7\baselineskip}
% mdparagraph:P114:begin
\section*{6. Practice in forms and word families}
% mdparagraph:P114:end

% mdparagraph:P115:begin
Give the forms or explanations requested. For analysis tasks use the supplied citation meanings and source forms. Their use here does not require you to produce unassigned forms from memory.
% mdparagraph:P115:end

\par\begin{minipage}{\linewidth}\raggedright
% mdparagraph:P116:begin
\subsection*{F01}
% mdparagraph:P116:end

% item: NEW-U14-F01
% mdparagraph:P117:begin
Give nominative, genitive, dative and accusative, singular and plural, of wǣdla and ēage. Explain the masculine/\allowbreak{}neuter contrast in the singular and which cells ending in -an require context to distinguish them.
% mdparagraph:P117:end

\end{minipage}\par\addvspace{4pt}

\par\begin{minipage}{\linewidth}\raggedright
% mdparagraph:P118:begin
\subsection*{F02}
% mdparagraph:P118:end

% item: NEW-U14-F02
% mdparagraph:P119:begin
Give the complete nominative, genitive, dative and accusative singular/\allowbreak{}plural of goldhord. Give the dative singular of wyrtgemang and gōdwębb. Explain where the case ending is attached in these compounds.
% mdparagraph:P119:end

\end{minipage}\par\addvspace{4pt}

\par\begin{minipage}{\linewidth}\raggedright
% mdparagraph:P120:begin
\subsection*{F03}
% mdparagraph:P120:end

% item: NEW-U14-F03
% mdparagraph:P121:begin
Give all six present and all six preterite indicative persons of liccian. Explain the present stems and the preterite stem shown by the paradigm.
% mdparagraph:P121:end

\end{minipage}\par\addvspace{4pt}

\par\begin{minipage}{\linewidth}\raggedright
% mdparagraph:P122:begin
\subsection*{F04}
% mdparagraph:P122:end

% item: NEW-U14-F04
% mdparagraph:P123:begin
Give these indicative forms: gewęndan present 2 singular; gelīefan present 3 singular; forgiefan present 3 singular; āhrēosan preterite 3 plural; gebiddan present 3 singular; ætīewan preterite 3 plural; āhębban present 3 singular; onfōn preterite 2 singular. Explain the stem choice in the two prefixed strong verbs āhrēosan and onfōn.
% mdparagraph:P123:end

\end{minipage}\par\addvspace{4pt}

\par\begin{minipage}{\linewidth}\raggedright
% mdparagraph:P124:begin
\subsection*{F05}
% mdparagraph:P124:end

% item: NEW-U14-F05
% mdparagraph:P125:begin
Give the complete present and preterite subjunctive of gebiddan. Then give present subjunctive 3 singular of forweorðan, 2 singular of sęndan, 3 plural of gelīefan and 3 singular of ārīsan.
% mdparagraph:P125:end

\end{minipage}\par\addvspace{4pt}

\par\begin{minipage}{\linewidth}\raggedright
% mdparagraph:P126:begin
\subsection*{F06}
% mdparagraph:P126:end

% item: NEW-U14-F06
% mdparagraph:P127:begin
Give present and past participle bases of forceorfan, weorpan, gemētan and geglęngan. Give singular and plural commands of gemiltsian and tilian.
% mdparagraph:P127:end

\end{minipage}\par\addvspace{4pt}

\par\begin{minipage}{\linewidth}\raggedright
% mdparagraph:P128:begin
\subsection*{A01}
% mdparagraph:P128:end

% item: NEW-U14-A01
% mdparagraph:P129:begin
Analyse goldhord and the source forms heofonfȳre and hēaclifum. For each, identify the two components and the final noun, separate any inflectional ending, and give the whole expression’s contextual meaning. Explain why a compound’s exact meaning must still be learned.
% mdparagraph:P129:end

% mdparagraph:P130:begin
\textbf{Inputs, in order:} goldhord · heofonfȳre · hēaclifum
% mdparagraph:P130:end

\end{minipage}\par\addvspace{4pt}

% teaching: U14-DESIGN-COMPOUND-BASES
% support: U14-DESIGN-COMPOUND-BASES
% mdparagraph:P131:begin
\textbf{Supplied help.} For this analysis, hord means treasure/\allowbreak{}hoard and clif means cliff. Heofonfȳre means fire from heaven (here lightning), dative singular; hēaclifum means high cliffs, dative plural. These source forms and component citations are supplied; do not reconstruct further forms from them.
% mdparagraph:P131:end

\par\begin{minipage}{\linewidth}\raggedright
% mdparagraph:P132:begin
\subsection*{A02}
% mdparagraph:P132:end

% item: NEW-U14-A02
% mdparagraph:P133:begin
Compare leornian/\allowbreak{}leornung, gefēra/\allowbreak{}gefērscipe, rihtwīs/\allowbreak{}rihtwīsnes and īdel/\allowbreak{}īdelnes. Identify the noun-forming suffix in each family, state the parts of speech and give the taught meanings. Explain why īdelnes cannot be translated merely by adding an English suffix to the previously taught meaning of īdel.
% mdparagraph:P133:end

% mdparagraph:P134:begin
\textbf{Inputs, in order:} leornian · leornung · gefēra · gefērscipe · rihtwīs · rihtwīsnes · īdel · īdelnes
% mdparagraph:P134:end

\end{minipage}\par\addvspace{4pt}

% teaching: U14-DESIGN-LEORNIAN-CITATION
% support: U14-DESIGN-LEORNIAN-CITATION
% mdparagraph:P135:begin
\textbf{Supplied help.} Leornian is the supplied infinitive citation of the already taught verb learn. This task analyses its family; it does not require independent infinitive recall or use.
% mdparagraph:P135:end

\par\begin{minipage}{\linewidth}\raggedright
% mdparagraph:P136:begin
\subsection*{A03}
% mdparagraph:P136:end

% item: NEW-U14-A03
% mdparagraph:P137:begin
Compare heofon/\allowbreak{}heofonlic, swīð/\allowbreak{}swīðe and sōð/\allowbreak{}sōðlice. Identify what forms an adjective and what forms an adverb; give each word’s supplied or taught meaning. Does the final e alone prove that a word is an adverb? Explain.
% mdparagraph:P137:end

% mdparagraph:P138:begin
\textbf{Inputs, in order:} heofon · heofonlic · swīð · swīðe · sōð · sōðlice
% mdparagraph:P138:end

\end{minipage}\par\addvspace{4pt}

% teaching: U14-DESIGN-ADJECTIVE-CITATIONS
% support: U14-DESIGN-ADJECTIVE-CITATIONS
% mdparagraph:P139:begin
\textbf{Supplied help.} Swīð strong and sōð true are supplied citation bases, as in the earlier adverb lesson. No complete adjective systems or independent use of these bases are assigned.
% mdparagraph:P139:end

\par\begin{minipage}{\linewidth}\raggedright
% mdparagraph:P140:begin
\subsection*{A04}
% mdparagraph:P140:end

% item: NEW-U14-A04
% mdparagraph:P141:begin
Compare rihtwīsnes/\allowbreak{}unrihtwīsnes, berstan/\allowbreak{}tōberstan and weorðan/\allowbreak{}forweorðan. Identify each added prefix and compare the taught whole-word meanings. Explain why the pairs do not license either free coinage or treating every prefixed verb as its unprefixed base.
% mdparagraph:P141:end

% mdparagraph:P142:begin
\textbf{Inputs, in order:} rihtwīsnes · unrihtwīsnes · berstan · tōberstan · weorðan · forweorðan
% mdparagraph:P142:end

\end{minipage}\par\addvspace{4pt}

% teaching: U14-DESIGN-PREFIX-CITATIONS
% support: U14-DESIGN-PREFIX-CITATIONS
% mdparagraph:P143:begin
\textbf{Supplied help.} The infinitive labels berstan, tōberstan, weorðan and forweorðan are supplied here as dictionary citations of already taught finite verbs. Analyse the labels; no independent infinitive recall or nonfinite use is required.
% mdparagraph:P143:end

\par\begin{minipage}{\linewidth}\raggedright
% mdparagraph:P144:begin
\subsection*{A05}
% mdparagraph:P144:end

% item: NEW-U14-A05
% mdparagraph:P145:begin
In the family rihtwīs → rihtwīsnes → rihtwīsnesse, distinguish derivation from inflection. Parse rihtwīsnesse in Ic lufie rihtwīsnesse. Would the isolated spelling rihtwīsnesse determine its case? Explain.
% mdparagraph:P145:end

% mdparagraph:P146:begin
\textbf{Inputs, in order:} Ic lufie rihtwīsnesse. · rihtwīs · rihtwīsnes
% mdparagraph:P146:end

\end{minipage}\par\addvspace{4pt}

\par\begin{minipage}{\linewidth}\raggedright
% mdparagraph:P147:begin
\subsection*{A06}
% mdparagraph:P147:end

% item: NEW-U14-A06
% mdparagraph:P148:begin
Translate Hē onfēng gōdum and Ic geseah ðā rihtwīsan. Parse gōdum as the assigned lexical noun and rihtwīsan as a substantival adjective. Explain each word’s case and the adjective’s agreement. Does the absence of an expressed following noun by itself settle a word’s lexical category?
% mdparagraph:P148:end

% mdparagraph:P149:begin
\textbf{Inputs, in order:} Hē onfēng gōdum. · Ic geseah ðā rihtwīsan.
% mdparagraph:P149:end

\end{minipage}\par\addvspace{4pt}

\reserve{7\baselineskip}
% mdparagraph:P150:begin
\section*{7. Short compositions}
% mdparagraph:P150:end

\par\begin{minipage}{\linewidth}\raggedright
% mdparagraph:P151:begin
\subsection*{C01}
% mdparagraph:P151:end

% item: NEW-U14-C01
% mdparagraph:P152:begin
Translate: The poor men saw the dogs and the wounds. Use wǣdla for poor man.
% mdparagraph:P152:end

\end{minipage}\par\addvspace{4pt}

\par\begin{minipage}{\linewidth}\raggedright
% mdparagraph:P153:begin
\subsection*{C02}
% mdparagraph:P153:end

% item: NEW-U14-C02
% mdparagraph:P154:begin
Translate: I leave you [plural] peace, and I give the poor man the treasure. Use forlǣtan, forgiefan, ðearfa and goldhord.
% mdparagraph:P154:end

\end{minipage}\par\addvspace{4pt}

\par\begin{minipage}{\linewidth}\raggedright
% mdparagraph:P155:begin
\subsection*{C03}
% mdparagraph:P155:end

% item: NEW-U14-C03
% mdparagraph:P156:begin
Translate: The poor men believe in God. Use wǣdla and gelīefan.
% mdparagraph:P156:end

\end{minipage}\par\addvspace{4pt}

\par\begin{minipage}{\linewidth}\raggedright
% mdparagraph:P157:begin
\subsection*{C04}
% mdparagraph:P157:end

% item: NEW-U14-C04
% mdparagraph:P158:begin
Translate as singular commands: Have mercy on the poor man, and send him the gold. Use gemiltsian, wǣdla and sęndan.
% mdparagraph:P158:end

\end{minipage}\par\addvspace{4pt}

\reserve{7\baselineskip}
% mdparagraph:P159:begin
\section*{8. Cook, Group VIII}
% mdparagraph:P159:end

% teaching: NEW-U14-LOCAL-HELPS
% mdparagraph:P160:begin
Translate the complete original prompts. Their parenthetical references are Cook’s grammar section numbers, not sections of this chapter. Modern directions and supplied helps follow separately. “Om.” means that the indicated English word has no separate Old English equivalent. A specified vocabulary or order target must be observed even when another grammatical expression is possible. Keep all help in view; the negative and indefinite expressions below are supplied for these items.
% mdparagraph:P160:end

\reserve{7\baselineskip}
% mdparagraph:P161:begin
\subsection*{Exercise 26}
% mdparagraph:P161:end

\par\begin{minipage}{\linewidth}\raggedright
% mdparagraph:P162:begin
\subsection*{Cook VIII.26.1}
% mdparagraph:P162:end

% item: EX-VIII-26-01
% mdparagraph:P163:begin
Thou art dust, and thou returnest (gewendan) to dust.
% mdparagraph:P163:end

\end{minipage}\par\addvspace{4pt}

\par\begin{minipage}{\linewidth}\raggedright
% mdparagraph:P164:begin
\subsection*{Cook VIII.26.2}
% mdparagraph:P164:end

% item: EX-VIII-26-02
% mdparagraph:P165:begin
So much (swīðe) loved God (om. the) world (middangeard), that he gave (sęllan) his only-begotten (āncęnned) Son, that no one (om. may) perish (196. f) that (87. d) believeth in (on, 172) him, but have (196. f) the eternal life.
% mdparagraph:P165:end

\end{minipage}\par\addvspace{4pt}

% teaching: U14-G-COOK-NONE
% support: U14-G-COOK-NONE
% mdparagraph:P166:begin
\textbf{Supplied help.} Nān means “no one” by itself, and “no” in nān mann. Also put ne with the finite verb; the two negatives express one negation.
% mdparagraph:P166:end

% teaching: U14-COOK-LEX-ANCENNED
% support: U14-COOK-LEX-ANCENNED
% mdparagraph:P167:begin
\textbf{Supplied help.} Use āncęnnedan, “only-begotten”: accusative masculine singular weak, agreeing with Sunu after his.
% mdparagraph:P167:end

% mdparagraph:P168:begin
\textbf{Directions.} The references ask for bare relative ðe, on + accusative with belief, and subjunctives in the final purpose. Omit a separate word for “may”.
% mdparagraph:P168:end

\par\begin{minipage}{\linewidth}\raggedright
% mdparagraph:P169:begin
\subsection*{Cook VIII.26.3}
% mdparagraph:P169:end

% item: EX-VIII-26-03
% mdparagraph:P170:begin
Jesus said (cweðan) to him, I am way, and truth, and life; no one (nān) cometh (183) to the Father, except by (ðurh) me.
% mdparagraph:P170:end

% support: U14-G-COOK-NONE
% mdparagraph:P171:begin
\textbf{Supplied help.} Nān means “no one” by itself, and “no” in nān mann. Also put ne with the finite verb; the two negatives express one negation.
% mdparagraph:P171:end

\end{minipage}\par\addvspace{4pt}

% teaching: U14-G-COOK-EXCEPT
% support: U14-G-COOK-EXCEPT
% mdparagraph:P172:begin
\textbf{Supplied help.} Use būtan, “except”, before ðurh + accusative. The supplied spelling ðurh corresponds to earlier þurh. Here būtan introduces the exception; it is not the earlier būtan + dative meaning “without”.
% mdparagraph:P172:end

% teaching: U14-G-COOK-SPEECH-TO
% support: U14-G-COOK-SPEECH-TO
% mdparagraph:P173:begin
\textbf{Supplied help.} Cwæð tō him means “said to him”. The speech addressee is dative after tō. Retain the first person within the direct quotation.
% mdparagraph:P173:end

% teaching: U14-COOK-LEX-WEG
% support: U14-COOK-LEX-WEG
% mdparagraph:P174:begin
\textbf{Supplied help.} Weg here means “way; means of access to the Father”, a figurative sense.
% mdparagraph:P174:end

\par\begin{minipage}{\linewidth}\raggedright
% mdparagraph:P175:begin
\subsection*{Cook VIII.26.4}
% mdparagraph:P175:end

% item: EX-VIII-26-04
% mdparagraph:P176:begin
Peace I leave (om. with) you (164. a); my peace I give (forgiefan) you.
% mdparagraph:P176:end

% mdparagraph:P177:begin
\textbf{Directions.} The recipients are plural and dative. Omit an Old English equivalent of English “with”; use the leave-peace construction in §5.
% mdparagraph:P177:end

\end{minipage}\par\addvspace{4pt}

\par\begin{minipage}{\linewidth}\raggedright
% mdparagraph:P178:begin
\subsection*{Cook VIII.26.5}
% mdparagraph:P178:end

% item: EX-VIII-26-05
% mdparagraph:P179:begin
Then rained it, and there (ðǣr) came floods, and blew winds, and fell (āhrēosan) on that house; and the house fell (feallan).
% mdparagraph:P179:end

\end{minipage}\par\addvspace{4pt}

% teaching: U14-G-COOK-RAIN
% support: U14-G-COOK-RAIN
% mdparagraph:P180:begin
\textbf{Supplied help.} Rīnde is the third singular preterite “rained”. Use impersonal hit, “it”, with this supplied form.
% mdparagraph:P180:end

\par\begin{minipage}{\linewidth}\raggedright
% mdparagraph:P181:begin
\subsection*{Cook VIII.26.6}
% mdparagraph:P181:end

% item: EX-VIII-26-06
% mdparagraph:P182:begin
So every good tree beareth (beran) good fruits (wǣstm), and every evil tree beareth evil fruits.
% mdparagraph:P182:end

% mdparagraph:P183:begin
\textbf{Directions.} C01: retain Cook’s wǣstm in the prompt, but write the earlier selected wæstm, with short æ. See the editorial note.
% mdparagraph:P183:end

\end{minipage}\par\addvspace{4pt}

\par\begin{minipage}{\linewidth}\raggedright
% mdparagraph:P184:begin
\subsection*{Cook VIII.26.7}
% mdparagraph:P184:end

% item: EX-VIII-26-07
% mdparagraph:P185:begin
Every tree that (87. d) beareth not good fruit (sing.) is cut down (forceorfan, 104) and cast into (om. the) fire (verb last).
% mdparagraph:P185:end

% mdparagraph:P186:begin
\textbf{Directions.} Use bare relative ðe. Keep “fruit” singular and omit the article before “fire”. For the stated order target, place the finite passive auxiliary last.
% mdparagraph:P186:end

\end{minipage}\par\addvspace{4pt}

\par\begin{minipage}{\linewidth}\raggedright
% mdparagraph:P187:begin
\subsection*{Cook VIII.26.8}
% mdparagraph:P187:end

% item: EX-VIII-26-08
% mdparagraph:P188:begin
He went into (on, 172) a mountain, that he (om. might) pray (gebiddan, 184. b); and when (ðā) he prayed (184. b), his raiment was white (hwīt), shining.
% mdparagraph:P188:end

\end{minipage}\par\addvspace{4pt}

% teaching: U14-G-COOK-SHINING
% support: U14-G-COOK-SHINING
% mdparagraph:P189:begin
\textbf{Supplied help.} Scīnende, “shining”, is a present participle agreeing with neuter singular rēaf. Keep it in the description of the raiment; no English-style progressive construction is required.
% mdparagraph:P189:end

% mdparagraph:P190:begin
\textbf{Directions.} Reference 184.b requires the accusative reflexive with both forms of gebiddan. The purpose clause and factual when-clause have different moods.
% mdparagraph:P190:end

\par\begin{minipage}{\linewidth}\raggedright
% mdparagraph:P191:begin
\subsection*{Cook VIII.26.9}
% mdparagraph:P191:end

% item: EX-VIII-26-09
% mdparagraph:P192:begin
The mountain, falling, (be)cometh to naught.
% mdparagraph:P192:end

\end{minipage}\par\addvspace{4pt}

% teaching: U14-G-COOK-FALL-NAUGHT
% support: U14-G-COOK-FALL-NAUGHT
% mdparagraph:P193:begin
\textbf{Supplied help.} Hrēosende, “falling”, is nominative masculine singular and describes munt. The fixed expression weorðan tō nāhte means “come to nothing”. Keep the supplied nāhte in this expression.
% mdparagraph:P193:end

\par\begin{minipage}{\linewidth}\raggedright
% mdparagraph:P194:begin
\subsection*{Cook VIII.26.10}
% mdparagraph:P194:end

% item: EX-VIII-26-10
% mdparagraph:P195:begin
The king, in all his glory (wuldor), was not adorned like (swā-swā) one of these.
% mdparagraph:P195:end

\end{minipage}\par\addvspace{4pt}

% teaching: U14-G-COOK-AS
% support: U14-G-COOK-AS
% mdparagraph:P196:begin
\textbf{Supplied help.} Use the connector swā-swā, “like; as”. The following ān + genitive ðissa means “one of these”.
% mdparagraph:P196:end

% teaching: U14-G-COOK-GLORY
% support: U14-G-COOK-GLORY
% mdparagraph:P197:begin
\textbf{Supplied help.} On eallum his wuldre means “in all his glory”. On governs the dative group; eallum agrees with neuter wuldre, while his identifies the possessor. Use ne wæs for “was not”.
% mdparagraph:P197:end

\reserve{7\baselineskip}
% mdparagraph:P198:begin
\subsection*{Exercise 27}
% mdparagraph:P198:end

\par\begin{minipage}{\linewidth}\raggedright
% mdparagraph:P199:begin
\subsection*{Cook VIII.27.1}
% mdparagraph:P199:end

% item: EX-VIII-27-01
% mdparagraph:P200:begin
He showed (ætīewan) them his precious things,—the silver, and the gold, and the spices, and all that was found (gemētan) in his treasures.
% mdparagraph:P200:end

% mdparagraph:P201:begin
\textbf{Directions.} The dative recipient differs from the accusative things shown. Strong adjective forms after a possessive remain possible, as taught earlier.
% mdparagraph:P201:end

\end{minipage}\par\addvspace{4pt}

\par\begin{minipage}{\linewidth}\raggedright
% mdparagraph:P202:begin
\subsection*{Cook VIII.27.2}
% mdparagraph:P202:end

% item: EX-VIII-27-02
% mdparagraph:P203:begin
A certain rich man was adorned (geglęngan) with purple (verb last).
% mdparagraph:P203:end

% mdparagraph:P204:begin
\textbf{Directions.} For the stated “verb last” target, put the finite auxiliary after both the material phrase and the participle.
% mdparagraph:P204:end

\end{minipage}\par\addvspace{4pt}

\par\begin{minipage}{\linewidth}\raggedright
% mdparagraph:P205:begin
\subsection*{Cook VIII.27.3}
% mdparagraph:P205:end

% item: EX-VIII-27-03
% mdparagraph:P206:begin
Then (ðā) lay a certain poor man at his gate.
% mdparagraph:P206:end

\end{minipage}\par\addvspace{4pt}

% teaching: U14-G-COOK-AT-GATES
% support: U14-G-COOK-AT-GATES
% mdparagraph:P207:begin
\textbf{Supplied help.} Use æt + dative place: æt his geate, “at his gate”; æt ēowrum gatum, “at your gates”. His names the possessor; ēowrum agrees with plural gatum.
% mdparagraph:P207:end

\par\begin{minipage}{\linewidth}\raggedright
% mdparagraph:P208:begin
\subsection*{Cook VIII.27.4}
% mdparagraph:P208:end

% item: EX-VIII-27-04
% mdparagraph:P209:begin
No man gave (om. unto) him; but dogs (hund, 43) came and licked (liccian, 118) his sores (wund, 43).
% mdparagraph:P209:end

% support: U14-G-COOK-NONE
% mdparagraph:P210:begin
\textbf{Supplied help.} Nān means “no one” by itself, and “no” in nān mann. Also put ne with the finite verb; the two negatives express one negation.
% mdparagraph:P210:end

% mdparagraph:P211:begin
\textbf{Directions.} C04: the noun wund, “wound; sore”, is feminine. Its complete table is in §2; the printed §43 reference is not its feminine declension. The dogs are plural subjects; the wounds belong to the poor man.
% mdparagraph:P211:end

\end{minipage}\par\addvspace{4pt}

\par\begin{minipage}{\linewidth}\raggedright
% mdparagraph:P212:begin
\subsection*{Cook VIII.27.5}
% mdparagraph:P212:end

% item: EX-VIII-27-05
% mdparagraph:P213:begin
Then happened it that the poor man departed (gewītan); and angels bore him to Abraham’s (54) dwelling.
% mdparagraph:P213:end

\end{minipage}\par\addvspace{4pt}

% teaching: U14-G-COOK-HAPPENED
% support: U14-G-COOK-HAPPENED
% mdparagraph:P214:begin
\textbf{Supplied help.} Gelamp means “happened”, singular past. A following ðæt-clause states what happened; optional introductory hit anticipates that clause. The event is factual, so use past indicative in the clause.
% mdparagraph:P214:end

% teaching: U14-COOK-LEX-ABRAHAM
% support: U14-COOK-LEX-ABRAHAM
% mdparagraph:P215:begin
\textbf{Supplied help.} Use Abrahames as the genitive of Abraham. Use unchanged Abraham as the accusative and in direct address.
% mdparagraph:P215:end

\par\begin{minipage}{\linewidth}\raggedright
% mdparagraph:P216:begin
\subsection*{Cook VIII.27.6}
% mdparagraph:P216:end

% item: EX-VIII-27-06
% mdparagraph:P217:begin
Then lifted (āhebban) he his eyes up, and saw Abraham afar off [from afar].
% mdparagraph:P217:end

\end{minipage}\par\addvspace{4pt}

% teaching: U14-G-COOK-AFAR
% support: U14-G-COOK-AFAR
% mdparagraph:P218:begin
\textbf{Supplied help.} Feorran means “from afar”. Its ending already expresses “from”; it needs no additional of.
% mdparagraph:P218:end

% support: U14-COOK-LEX-ABRAHAM
% mdparagraph:P219:begin
\textbf{Supplied help.} Use Abrahames as the genitive of Abraham. Use unchanged Abraham as the accusative and in direct address.
% mdparagraph:P219:end

\par\begin{minipage}{\linewidth}\raggedright
% mdparagraph:P220:begin
\subsection*{Cook VIII.27.7}
% mdparagraph:P220:end

% item: EX-VIII-27-07
% mdparagraph:P221:begin
O father Abraham, have compassion on me, and send Lazarus.
% mdparagraph:P221:end

% support: U14-COOK-LEX-ABRAHAM
% mdparagraph:P222:begin
\textbf{Supplied help.} Use Abrahames as the genitive of Abraham. Use unchanged Abraham as the accusative and in direct address.
% mdparagraph:P222:end

\end{minipage}\par\addvspace{4pt}

% teaching: U14-COOK-LEX-LAZARUS
% support: U14-COOK-LEX-LAZARUS
% mdparagraph:P223:begin
\textbf{Supplied help.} Use Lazarum as accusative singular. Lazaras is accusative plural, “Lazaruses”: people like Lazarus.
% mdparagraph:P223:end

\par\begin{minipage}{\linewidth}\raggedright
% mdparagraph:P224:begin
\subsection*{Cook VIII.27.8}
% mdparagraph:P224:end

% item: EX-VIII-27-08
% mdparagraph:P225:begin
O son, remember that thou didst receive goods in thy life.
% mdparagraph:P225:end

\end{minipage}\par\addvspace{4pt}

% teaching: U14-G-COOK-REMEMBER
% support: U14-G-COOK-REMEMBER
% mdparagraph:P226:begin
\textbf{Supplied help.} Gemun is the singular command “remember”. Its following ðæt-clause recalls an asserted past fact; use indicative onfēnge for “thou didst receive”.
% mdparagraph:P226:end

% teaching: U14-G-COOK-LIFETIME
% support: U14-G-COOK-LIFETIME
% mdparagraph:P227:begin
\textbf{Supplied help.} On ðīnum līfe means “in; during your life”. The whole lifetime phrase is dative singular.
% mdparagraph:P227:end

\par\begin{minipage}{\linewidth}\raggedright
% mdparagraph:P228:begin
\subsection*{Cook VIII.27.9}
% mdparagraph:P228:end

% item: EX-VIII-27-09
% mdparagraph:P229:begin
I beseech thee that thou send him to my father’s house.
% mdparagraph:P229:end

% mdparagraph:P230:begin
\textbf{Directions.} Use the earlier hālsian + ðæt request frame. Distinguish the addressee from the person sent and keep the internal genitive in “my father’s house”.
% mdparagraph:P230:end

\end{minipage}\par\addvspace{4pt}

\par\begin{minipage}{\linewidth}\raggedright
% mdparagraph:P231:begin
\subsection*{Cook VIII.27.10}
% mdparagraph:P231:end

% item: EX-VIII-27-10
% mdparagraph:P232:begin
If they hear not the prophets, they (om. will) not believe, though any one (hwā) arise from death (verb last).
% mdparagraph:P232:end

\end{minipage}\par\addvspace{4pt}

% teaching: U14-G-COOK-INDEFINITE-CONCESSIVE
% support: U14-G-COOK-INDEFINITE-CONCESSIVE
% mdparagraph:P233:begin
\textbf{Supplied help.} Hwā here has the indefinite meaning “anyone”. In the concessive clause use present subjunctive ārīse and obey the requested final-verb order.
% mdparagraph:P233:end

% mdparagraph:P234:begin
\textbf{Directions.} The present belief form can have future reference here; omit an added “will”. The last-verb instruction applies inside the concessive clause.
% mdparagraph:P234:end

\par\begin{minipage}{\linewidth}\raggedright
% mdparagraph:P235:begin
\subsection*{Cook VIII.27.11}
% mdparagraph:P235:end

% item: EX-VIII-27-11
% mdparagraph:P236:begin
Provide for yourselves (184. a) friends (156. n) in (on) God’s poor (ðearfa). Ye have many Lazaruses (Lazaras) now lying at your gates.
% mdparagraph:P236:end

\end{minipage}\par\addvspace{4pt}

% teaching: U14-G-COOK-LYING
% support: U14-G-COOK-LYING
% mdparagraph:P237:begin
\textbf{Supplied help.} Licgende, “lying”, is accusative masculine plural, describing the people denoted by Lazaras. Habbað has these people as its object; the participle describes their situation at the gates.
% mdparagraph:P237:end

% teaching: U14-G-COOK-POOR
% support: U14-G-COOK-POOR
% mdparagraph:P238:begin
\textbf{Supplied help.} In this expression on Godes ðearfum means “in; among God’s poor”. The group after on is dative plural; Godes is its dependent genitive.
% mdparagraph:P238:end

% support: U14-G-COOK-AT-GATES
% mdparagraph:P239:begin
\textbf{Supplied help.} Use æt + dative place: æt his geate, “at his gate”; æt ēowrum gatum, “at your gates”. His names the possessor; ēowrum agrees with plural gatum.
% mdparagraph:P239:end

% support: U14-COOK-LEX-LAZARUS
% mdparagraph:P240:begin
\textbf{Supplied help.} Use Lazarum as accusative singular. Lazaras is accusative plural, “Lazaruses”: people like Lazarus.
% mdparagraph:P240:end

% mdparagraph:P241:begin
\textbf{Directions.} Reference 184.a selects the dative reflexive beneficiary; 156.n selects the genitive object of tilian. The second sentence still addresses more than one person.
% mdparagraph:P241:end

\reserve{7\baselineskip}
% mdparagraph:P242:begin
\section*{9. Reading: The Dangers of Greatness}
% mdparagraph:P242:end

% teaching: NEW-U14-SOURCE-READING-SUPPORT
% teaching: NEW-U14-SOURCE-HELPS
% teaching: NEW-U14-GRAM-U14-SOURCE-ARGUMENT-STRUCTURE
% teaching: NEW-U14-GRAM-U14-LEXICAL-PARALLELISM
% mdparagraph:P243:begin
The following authentic passage develops a general claim through concrete comparisons and returns to human affairs. Read it as one argument. Repeated words, paired verbs and parallel clauses can reinforce one idea rather than add separate events. Keep track of each subject, its dependent clauses and any pronoun referring back to it.
% mdparagraph:P243:end

% mdparagraph:P244:begin
Cook heads the reading “From Wulfstan’s Homilies, No. 49”, but his introduction questions the attribution. His full qualification is retained in the source notes. The assigned opening ends with tō nāhte; the following argument is not assigned. Page and prose-line divisions and apparatus cues are removed from the reading paragraph below. The linguistic spelling, including source hyphens and vowel marks, is retained; the single page division within geslægene is joined. Note attachments are listed after the reading questions.\textsuperscript{6}
% mdparagraph:P244:end

\par\begin{minipage}{\linewidth}\raggedright
% source: FB-VIII-DANGERS-OPENING-ARGUMENT
% mdparagraph:P245:begin
\readingtext{Ǣghwilc hēah ār, hēr on worulde, bið mid frēcnessum ymbseald; efne swā ðā woruldgeðyngða bēoð māran, swā ðā frēcnessa bēoð swīðran. Swā wē magon, be ðǣm, ðā bȳsena oncnāwan and ongietan. Ðæt trēow, ðonne, ðe wiexð on ðǣm wudubearwe, ðæt hit hlīfað ūp ofer eall ðā ōðru trēowu and brǣt hit, ðonne sęmninga storm gestęnt, and se stranga wind, ðonne bið hit swīðlicor gewǣged and geswęnged ðonne se ōðer wudu. Swā bið ēac gelīce be ðǣm hēaclifum and torrum, ðonne hīe hlīfiað feorr ūp ofer ðā ōðre eorðan, hīe ðonne sęmninga feallan onginnað, and full ðearlice hrēosan tō eorðan. Swilce ēac be ðǣm hēagum muntum and dūnum, ðā-ðe hēah standað ofer ealne middangeard, ðā-hwæðre wīte habbað ðæs ealdordōmes, ðæt hīe bēoð geneahhe mid heofonfȳre geðrēade and geðrǣste, and mid līegum geslægene. Swā ðā hēan mihta hēr on worulde hrēosað, and feallað, and tō lore weorðað, and ðisse worulde welan weorðað tō sorge, and ðās eorðlican wundor weorðað tō nāhte.}
% mdparagraph:P245:end

\end{minipage}\par
\reserve{7\baselineskip}
% mdparagraph:P246:begin
\subsection*{Vocabulary supplied for this reading}
% mdparagraph:P246:end

% mdparagraph:P247:begin
These meanings and source-specific forms remain visible for S01–S07. They do not introduce complete new paradigms. Nouns and adjectives not marked as newly productive in §2 are interpreted here through this help.
% mdparagraph:P247:end

% teaching: U14-READ-LEX-AEGHWILC
% support: U14-READ-LEX-AEGHWILC
% mdparagraph:P248:begin
Ǣghwilc — every; each (nominative feminine singular, modifying ār).
% mdparagraph:P248:end

% teaching: U14-READ-LEX-AR
% support: U14-READ-LEX-AR
% mdparagraph:P249:begin
ār — honour; high office or position (nominative feminine singular; subject).
% mdparagraph:P249:end

% teaching: U14-READ-LEX-FRECNES
% support: U14-READ-LEX-FRECNES
% mdparagraph:P250:begin
frēcnessum — danger; peril (dative feminine plural after mid); frēcnessa — danger; peril (nominative feminine plural; subject).
% mdparagraph:P250:end

% teaching: U14-READ-LEX-WORULDGETHYNGTH
% support: U14-READ-LEX-WORULDGETHYNGTH
% mdparagraph:P251:begin
woruldgeðyngða — worldly honour; worldly dignity (nominative feminine plural; subject).
% mdparagraph:P251:end

% teaching: U14-READ-LEX-BYSEN
% support: U14-READ-LEX-BYSEN
% mdparagraph:P252:begin
bȳsena — example; illustration (accusative feminine plural; object of oncnāwan and ongietan).
% mdparagraph:P252:end

% teaching: U14-READ-LEX-WUDUBEARU
% support: U14-READ-LEX-WUDUBEARU
% mdparagraph:P253:begin
wudubearwe — forest; grove (dative masculine singular after on).
% mdparagraph:P253:end

% teaching: U14-READ-LEX-WUDU
% support: U14-READ-LEX-WUDU
% mdparagraph:P254:begin
wudu — wood; forest, collectively (nominative masculine singular in comparative subject phrase).
% mdparagraph:P254:end

% teaching: U14-READ-LEX-HEACLIF
% support: U14-READ-LEX-HEACLIF
% mdparagraph:P255:begin
hēaclifum — high cliff (dative neuter plural after be; transparent hēah + clif compound).
% mdparagraph:P255:end

% teaching: U14-READ-LEX-DUN
% support: U14-READ-LEX-DUN
% mdparagraph:P256:begin
dūnum — mountain; hill (dative feminine plural after be).
% mdparagraph:P256:end

% teaching: U14-READ-LEX-WITE
% support: U14-READ-LEX-WITE
% mdparagraph:P257:begin
wīte — punishment; penalty (accusative neuter singular; object of habbað).
% mdparagraph:P257:end

% teaching: U14-READ-LEX-HEOFONFYR
% support: U14-READ-LEX-HEOFONFYR
% mdparagraph:P258:begin
heofonfȳre — fire from heaven; heavenly fire (dative neuter singular of instrument after mid).
% mdparagraph:P258:end

% teaching: U14-READ-LEX-LIEG
% support: U14-READ-LEX-LIEG
% mdparagraph:P259:begin
līegum — thunderbolt; levin (lightning) (dative masculine plural of instrument after mid).
% mdparagraph:P259:end

% teaching: U14-READ-LEX-LOR
% support: U14-READ-LEX-LOR
% mdparagraph:P260:begin
lore — destruction; ruin (dative neuter singular in supplied tō lore weorðan).
% mdparagraph:P260:end

% teaching: U14-READ-LEX-WELA
% support: U14-READ-LEX-WELA
% mdparagraph:P261:begin
welan — riches; wealth (nominative masculine plural; subject; modified by genitive ðisse worulde).
% mdparagraph:P261:end

% teaching: U14-READ-LEX-SORG
% support: U14-READ-LEX-SORG
% mdparagraph:P262:begin
sorge — distress; anxiety; trouble; sorrow (dative feminine singular in supplied result tō sorge).
% mdparagraph:P262:end

% teaching: U14-READ-LEX-EORTHLIC
% support: U14-READ-LEX-EORTHLIC
% mdparagraph:P263:begin
eorðlican — earthly (nominative neuter plural weak, agreeing with wundor after ðās).
% mdparagraph:P263:end

% teaching: U14-READ-LEX-TREO
% support: U14-READ-LEX-TREO
% mdparagraph:P264:begin
trēow — tree (nominative neuter singular, source -w variant of taught trēo).
% mdparagraph:P264:end

% teaching: U14-READ-LEX-HEAH
% support: U14-READ-LEX-HEAH
% mdparagraph:P265:begin
hēah — high in rank; exalted (nominative feminine singular exact source form, modifying ār; beside earlier feminine hēa); hēagum — high (dative plural weak after ðǣm, modifying coordinated muntum and dūnum; the source retains g); hēan — high in rank; exalted.
% mdparagraph:P265:end

% teaching: U14-READ-LEX-SWITHRA
% support: U14-READ-LEX-SWITHRA
% mdparagraph:P266:begin
swīðran — stronger; fiercer (degree), distinct from earlier right-side sense.
% mdparagraph:P266:end

% teaching: U14-READ-LEX-STORM
% support: U14-READ-LEX-STORM
% mdparagraph:P267:begin
storm — storm (nominative masculine singular, first subject of gestęnt).
% mdparagraph:P267:end

% teaching: U14-READ-LEX-STRANG
% support: U14-READ-LEX-STRANG
% mdparagraph:P268:begin
stranga — strong, here attributive with wind (nominative masculine singular weak after se).
% mdparagraph:P268:end

% teaching: U14-READ-LEX-OTHER
% support: U14-READ-LEX-OTHER
% mdparagraph:P269:begin
ōðer — remaining; the rest of the (forest/\allowbreak{}earth); ōðre — remaining; the rest of the (forest/\allowbreak{}earth).
% mdparagraph:P269:end

% teaching: U14-READ-LEX-TORR
% support: U14-READ-LEX-TORR
% mdparagraph:P270:begin
torrum — crag, the probable sense here in Cook’s note.
% mdparagraph:P270:end

% teaching: U14-READ-LEX-EALDORDOM
% support: U14-READ-LEX-EALDORDOM
% mdparagraph:P271:begin
ealdordōmes — chief place; primacy; eminence (genitive masculine singular dependent on wīte).
% mdparagraph:P271:end

% teaching: U14-READ-LEX-MIHT
% support: U14-READ-LEX-MIHT
% mdparagraph:P272:begin
mihta — powers; those holding high power or station in the concluding human application (nominative feminine plural; subject of hrēosað, feallað and weorðað).
% mdparagraph:P272:end

\reserve{7\baselineskip}
% mdparagraph:P273:begin
\subsection*{Supplied clause help}
% mdparagraph:P273:end

% teaching: U14-G-SRC-01
% support: U14-G-SRC-01
% mdparagraph:P274:begin
\textbf{Clause help.} Ymbseald means “surrounded; beset”; mid introduces the dangers. Efne emphasizes the statement. Swā … māran, swā … swīðran means “the greater …, the fiercer …”. This supplied proportional comparison differs from equal-degree swā … swā.
% mdparagraph:P274:end

% teaching: U14-G-SRC-02
% support: U14-G-SRC-02
% mdparagraph:P275:begin
\textbf{Clause help.} Be ðǣm is a separate dative phrase, “by; from that”, pointing back to the preceding claim. Ðā bȳsena is the accusative object of both oncnāwan and ongietan, infinitives after magon. The paired verbs reinforce one act of understanding.
% mdparagraph:P275:end

% teaching: U14-G-SRC-03
% support: U14-G-SRC-03
% mdparagraph:P276:begin
\textbf{Clause help.} Relative ðe identifies the tree. Wiexð is the source form “grows”, beside earlier wiext. The result clause ðæt hit hlīfað means “so that it towers”; this hit is nominative. Brǣt hit means “spreads itself”: this hit is accusative and refers to the same tree.
% mdparagraph:P276:end

% teaching: U14-G-SRC-04
% support: U14-G-SRC-04
% mdparagraph:P277:begin
\textbf{Clause help.} Sęmninga means “suddenly”; gestęnt, “assails; comes on”. Storm and se stranga wind are its coordinated subjects, although the source retains singular gestęnt. Gewǣged and geswęnged mean “shaken; buffeted”, passive participles. The main-clause ðonne is “then”; the ðonne following swīðlicor, “more violently”, is “than”. Se ōðer wudu means “the rest of the forest”.
% mdparagraph:P277:end

% teaching: U14-G-SRC-HREOSAN-NOTE
% support: U14-G-SRC-HREOSAN-NOTE
% mdparagraph:P278:begin
\textbf{Clause help.} For the question in original note 157.3, the indicative preterite plural of hrēosan is hruron, in all three persons. The plural stem has u and r in place of present ēo and s. This answer is supplied reference, not an additional exercise.
% mdparagraph:P278:end

% teaching: U14-G-SRC-05
% support: U14-G-SRC-05
% mdparagraph:P279:begin
\textbf{Clause help.} Swā bið ēac gelīce be … means “it is likewise also with; concerning …”. Hīe refers to the cliffs and crags. Hlīfiað means “tower”; feorr ūp, “far up”. Onginnað, “begin”, governs both feallan and hrēosan. Full ðearlice means “very violently”; the paired verbs describe one collapse. Tō eorðan gives its destination.
% mdparagraph:P279:end

% teaching: U14-G-SRC-RELATIVE
% support: U14-G-SRC-RELATIVE
% mdparagraph:P280:begin
\textbf{Clause help.} Ðā-ðe is the source’s joined spelling of the plural relative ðā + ðe, “which”. It is nominative plural, subject of standað, and refers to the mountains named by muntum and dūnum.
% mdparagraph:P280:end

% teaching: U14-G-SRC-06
% support: U14-G-SRC-06
% mdparagraph:P281:begin
\textbf{Clause help.} After Swilce ēac supply “it is likewise with; concerning”. Hēah is adverbial “high”; ðā-hwæðre means “nevertheless”; geneahhe, “often”. The ðæt-clause explains the penalty of pre-eminence. Geðrēade, geðrǣste and geslægene are agreeing plural passive participles, “punished; battered; struck”, with heavenly fire and thunderbolts as the instruments.
% mdparagraph:P281:end

% teaching: U14-G-SRC-07
% support: U14-G-SRC-07
% mdparagraph:P282:begin
\textbf{Clause help.} The high powers hrēosað and feallað, two verbs reinforcing their fall. Tō lore weorðan means “come to destruction”; tō sorge, “to sorrow”; tō nāhte, “to nothing”. These fixed tō-result expressions differ from weorðan followed by a nominative predicate.
% mdparagraph:P282:end

% mdparagraph:P283:begin
For the original question in note 157.6, modern English a down is a hill or upland. The adverb down developed from an expression meaning “from a hill”. This historical information is supplied; no additional etymology exercise is assigned.
% mdparagraph:P283:end

\reserve{7\baselineskip}
% mdparagraph:P284:begin
\subsection*{Reading questions}
% mdparagraph:P284:end

\par\begin{minipage}{\linewidth}\raggedright
% mdparagraph:P285:begin
\subsection*{S01}
% mdparagraph:P285:end

% item: NEW-U14-S01
% mdparagraph:P286:begin
Translate the complete assigned Dangers of Greatness passage, using its local helps and notes. Preserve each comparison and the final three outcomes.
% mdparagraph:P286:end

\end{minipage}\par\addvspace{4pt}

\par\begin{minipage}{\linewidth}\raggedright
% mdparagraph:P287:begin
\subsection*{S02}
% mdparagraph:P287:end

% item: NEW-U14-S02
% mdparagraph:P288:begin
State the opening claim, identify the three examples in their order, and explain how the final sentence applies them to human affairs.
% mdparagraph:P288:end

\end{minipage}\par\addvspace{4pt}

\par\begin{minipage}{\linewidth}\raggedright
% mdparagraph:P289:begin
\subsection*{S03}
% mdparagraph:P289:end

% item: NEW-U14-S03
% mdparagraph:P290:begin
In the tree sentence, identify the referent of each hit, explain the relative and result clauses, distinguish the temporal ðonne … ðonne from the final comparative ðonne, and explain se ōðer wudu.
% mdparagraph:P290:end

\end{minipage}\par\addvspace{4pt}

\par\begin{minipage}{\linewidth}\raggedright
% mdparagraph:P291:begin
\subsection*{S04}
% mdparagraph:P291:end

% item: NEW-U14-S04
% mdparagraph:P292:begin
In the cliff sentence, identify what hīe refers to and which verb governs feallan and hrēosan. Explain how the infinitives and the adverbial phrase full ðearlice contribute to the single event.
% mdparagraph:P292:end

\end{minipage}\par\addvspace{4pt}

\par\begin{minipage}{\linewidth}\raggedright
% mdparagraph:P293:begin
\subsection*{S05}
% mdparagraph:P293:end

% item: NEW-U14-S05
% mdparagraph:P294:begin
Analyse ealdordōmes and eorðlican in the source. Separate the derivational element from the inflection, identify each word’s part of speech and contextual role, and explain the whole expression in English. Use the supplied source analyses; no additional forms are required.
% mdparagraph:P294:end

\end{minipage}\par\addvspace{4pt}

% teaching: U14-DESIGN-SOURCE-DERIVATIONS
% support: U14-DESIGN-SOURCE-DERIVATIONS
% mdparagraph:P295:begin
\textbf{Supplied help.} Ealdor is supplied here as chief/\allowbreak{}king; ealdordōm is supremacy or pre-eminence. The source ealdordōmes is genitive singular, dependent on wīte punishment/\allowbreak{}penalty. Eorðlic is earthly; eorðlican is weak nominative neuter plural agreeing with ðās … wundor. These exact contextual analyses and citation bases are supplied for formation analysis, not assigned as complete new paradigms.
% mdparagraph:P295:end

\par\begin{minipage}{\linewidth}\raggedright
% mdparagraph:P296:begin
\subsection*{S06}
% mdparagraph:P296:end

% item: NEW-U14-S06
% mdparagraph:P297:begin
In the final sentence, list the three subjects with their stated outcomes. Explain why ðisse worulde is genitive while ðās agrees with wundor.
% mdparagraph:P297:end

\end{minipage}\par\addvspace{4pt}

\par\begin{minipage}{\linewidth}\raggedright
% mdparagraph:P298:begin
\subsection*{S07}
% mdparagraph:P298:end

% item: NEW-U14-S07
% mdparagraph:P299:begin
Choose two pairs or repetitions in the passage that reinforce its argument. Quote the Old English words and explain their effect in context. Do not assume that paired near-synonyms must describe separate events.
% mdparagraph:P299:end

\end{minipage}\par\addvspace{4pt}

\reserve{8\baselineskip}
% mdparagraph:P300:begin
\section*{10. Original source apparatus}
% mdparagraph:P300:end

\begingroup\fontsize{10}{12.8}\selectfont
\setlength{\parskip}{1.5pt}
% mdparagraph:P301:begin
These are Cook’s attached notes, including his questions and literary parallels. They are historical apparatus, not extra assigned exercises. Modern clause help above supplies the needed grammar. Note labels combine the printed page and original note number; prose-line positions in the attachment table are editorial counts anchored by Cook’s printed marginal numbers.
% mdparagraph:P301:end

\reserve{7\baselineskip}
% mdparagraph:P302:begin
\subsection*{Cook’s introductory notice}
% mdparagraph:P302:end

% mdparagraph:P303:begin
VIII. THE DANGERS OF GREATNESS.
% mdparagraph:P303:end

% mdparagraph:P304:begin
(From Wulfstan’s Homilies, No. 49.)
% mdparagraph:P304:end

% mdparagraph:P305:begin
[Wulfstan — also known by his Latinized name, Lupus — was Bishop of Worcester and Archbishop of York from 1002 to 1023. This homily is one of those attributed to him, but, according to Napier, with insufficient reason, as a portion of it is found in the Blickling Homilies, the manuscript of which bears the date of 971.]
% mdparagraph:P305:end

\reserve{7\baselineskip}
% mdparagraph:P306:begin
\subsection*{Note attachments}
% mdparagraph:P306:end

% mdtable:T43:begin
\par\begingroup
\fontsize{10.5}{13.2}\selectfont
\setlength{\LTleft}{0pt}\setlength{\LTright}{\fill}
\begin{longtable}{@{}>{\raggedright\arraybackslash}p{24.16mm}>{\raggedright\arraybackslash}p{24.16mm}>{\raggedright\arraybackslash}p{70.06mm}@{}}
\toprule
\textbf{Note} & \textbf{Page.line} & \textbf{Immediately after this source text} \\
\midrule
\endfirsthead
\toprule
\textbf{Note} & \textbf{Page.line} & \textbf{Immediately after this source text} \\
\midrule
\endhead
\midrule
\multicolumn{3}{r}{\footnotesize Continued on the following page.}\\
\endfoot
\bottomrule
\endlastfoot
156.1 & 156.1 & Ǣghwilc hēah ār, hēr on worulde, bið mid frēcnessum \\
156.2 & 156.2 & ymbseald \\
156.3 & 156.2 & ymbseald; efne swā \\
156.4 & 156.4 & bȳsena oncnāwan and ongietan. \\
156.5 & 156.5 & wiexð \\
156.6 & 156.5 & wiexð on ðǣm wudubearwe, ðæt \\
156.6a & 156.6 & ðā ōðru trēowu and brǣt \\
156.7 & 156.6 & ðā ōðru trēowu and brǣt hit, \\
156.8 & 156.6 & ðā ōðru trēowu and brǣt hit, ðonne sęmninga storm \\
156.9 & 156.7 & gestęnt, and se stranga wind, \\
156.10 & 156.7 & gestęnt, and se stranga wind, ðonne \\
156.11 & 156.8 & gewǣged and geswęnged ðonne se ōðer wudu. \\
156.12 & 156.9 & ēac gelīce be ðǣm hēaclifum and torrum, \\
157.1 & 157.1 & hlīfiað feorr ūp ofer ðā ōðre \\
157.2 & 157.2 & feallan onginnað, \\
157.3 & 157.2 & feallan onginnað, and full ðearlice hrēosan \\
157.4 & 157.3 & Swilce \\
157.5 & 157.3 & Swilce ēac be ðǣm hēagum \\
157.6 & 157.3 & Swilce ēac be ðǣm hēagum muntum and dūnum, \\
157.7 & 157.6 & heofonfȳre \\
158.1 & 158.1 & slægene. Swā ðā hēan mihta \\
158.2 & 158.2 & and feallað, \\
158.3 & 158.2 & and feallað, and tō lore weorðað, and ðisse \\
158.3 & 158.2 & and feallað, and tō lore weorðað, and ðisse worulde \\
158.4 & 158.4 & weorðað tō nāhte. \\
\end{longtable}
\par\endgroup

% mdtable:T43:end
\reserve{7\baselineskip}
% mdparagraph:P307:begin
\subsection*{156.1}
% mdparagraph:P307:end

% source-note: FB-VIII-DANGERS-156-N1
% mdparagraph:P308:begin
See 144.
% mdparagraph:P308:end

\reserve{7\baselineskip}
% mdparagraph:P309:begin
\subsection*{156.2}
% mdparagraph:P309:end

% source-note: FB-VIII-DANGERS-156-N2
% mdparagraph:P310:begin
See 114.
% mdparagraph:P310:end

\reserve{7\baselineskip}
% mdparagraph:P311:begin
\subsection*{156.3}
% mdparagraph:P311:end

% source-note: FB-VIII-DANGERS-156-N3
% mdparagraph:P312:begin
Swā . . . māran, swā . . . swīðran = the greater, the fiercer. Note the tendency to antithesis.
% mdparagraph:P312:end

\reserve{7\baselineskip}
% mdparagraph:P313:begin
\subsection*{156.4}
% mdparagraph:P313:end

% source-note: FB-VIII-DANGERS-156-N4
% mdparagraph:P314:begin
Observe the redundancy.
% mdparagraph:P314:end

\reserve{7\baselineskip}
% mdparagraph:P315:begin
\subsection*{156.5}
% mdparagraph:P315:end

% source-note: FB-VIII-DANGERS-156-N5
% mdparagraph:P316:begin
See weaxan.
% mdparagraph:P316:end

\reserve{7\baselineskip}
% mdparagraph:P317:begin
\subsection*{156.6}
% mdparagraph:P317:end

% source-note: FB-VIII-DANGERS-156-N6
% mdparagraph:P318:begin
= so that.
% mdparagraph:P318:end

\reserve{7\baselineskip}
% mdparagraph:P319:begin
\subsection*{156.6a}
% mdparagraph:P319:end

% source-note: FB-VIII-DANGERS-156-N6a
% mdparagraph:P320:begin
See brǣdan (34).
% mdparagraph:P320:end

\reserve{7\baselineskip}
% mdparagraph:P321:begin
\subsection*{156.7}
% mdparagraph:P321:end

% source-note: FB-VIII-DANGERS-156-N7
% mdparagraph:P322:begin
See 184. b.
% mdparagraph:P322:end

\reserve{7\baselineskip}
% mdparagraph:P323:begin
\subsection*{156.8}
% mdparagraph:P323:end

% source-note: FB-VIII-DANGERS-156-N8
% mdparagraph:P324:begin
Note the alliteration.
% mdparagraph:P324:end

\reserve{7\baselineskip}
% mdparagraph:P325:begin
\subsection*{156.9}
% mdparagraph:P325:end

% source-note: FB-VIII-DANGERS-156-N9
% mdparagraph:P326:begin
Second subject of gestęnt.
% mdparagraph:P326:end

\reserve{7\baselineskip}
% mdparagraph:P327:begin
\subsection*{156.10}
% mdparagraph:P327:end

% source-note: FB-VIII-DANGERS-156-N10
% mdparagraph:P328:begin
Frequently the second correlative, in such pairs as ðonne . . . ðonne, ðā . . . ðā, need not be translated; it is frequently followed by an inverted order, as here, the verb preceding its subject. See 202.
% mdparagraph:P328:end

\reserve{7\baselineskip}
% mdparagraph:P329:begin
\subsection*{156.11}
% mdparagraph:P329:end

% source-note: FB-VIII-DANGERS-156-N11
% mdparagraph:P330:begin
Se ōðer wudu = the rest of the forest, not the other wood.
% mdparagraph:P330:end

\reserve{7\baselineskip}
% mdparagraph:P331:begin
\subsection*{156.12}
% mdparagraph:P331:end

% source-note: FB-VIII-DANGERS-156-N12
% mdparagraph:P332:begin
Probably here = crag.
% mdparagraph:P332:end

\reserve{7\baselineskip}
% mdparagraph:P333:begin
\subsection*{157.1}
% mdparagraph:P333:end

% source-note: FB-VIII-DANGERS-157-N1
% mdparagraph:P334:begin
See p. 156, n. 12.
% mdparagraph:P334:end

\reserve{7\baselineskip}
% mdparagraph:P335:begin
\subsection*{157.2}
% mdparagraph:P335:end

% source-note: FB-VIII-DANGERS-157-N2
% mdparagraph:P336:begin
This resembles the use of gin in Chaucer, almost as an auxiliary tense-sign, like do in Mod. Eng., the latter not being thus used in OE. In Chaucer it usually occurs as the preterit gan, e.g. in the Clerk’s Tale, 392: “til the sonne gan descende.” See Lounsbury’s History of the English Language. An interesting parallel is to be found in New Testament Greek, as, for example, Acts 1. 1: “The former treatise I made, O Theophilus, concerning all that Jesus began (ἤρξατο) both to do and to teach.” According to Thayer, however (Greek-English Lexicon of the New Testament), there is in its employment always a sense of beginning, in its proper meaning.
% mdparagraph:P336:end

\reserve{7\baselineskip}
% mdparagraph:P337:begin
\subsection*{157.3}
% mdparagraph:P337:end

% source-note: FB-VIII-DANGERS-157-N3
% mdparagraph:P338:begin
Dependent, like feallan, on onginnað. Give the ind. pret. plur.
% mdparagraph:P338:end

\reserve{7\baselineskip}
% mdparagraph:P339:begin
\subsection*{157.4}
% mdparagraph:P339:end

% source-note: FB-VIII-DANGERS-157-N4
% mdparagraph:P340:begin
It would be interesting to know from what literary source these illustrations are ultimately derived. They remind one of Shakespeare (Rich. III., 1. 3. 259–260):— They that stand high have many blasts to shake them; And if they fall, they dash themselves to pieces. Cf. also 3 Hen. VI., 5. 2. 11–15. No doubt many Elizabethan parallels could be found; I have noted in Chapman, Byron’s Conspiracy, Act 3, Scene 1 (p. 232 of Shepherd’s ed.), and Byron’s Tragedy, Act 5, Scene 1 (Ib., p. 272). Perhaps the Elizabethans may have derived them from Seneca; cf. the Chorus in Act 4 of the Hippolytus, vv. 1123–1143; Hercules Furens 201; Œdipus 8–11. Seneca may have caught a suggestion from Sophocles, though the parallel is somewhat remote; see the latter’s Antigone, vv. 712–717, and Horace, C. II. 10.
% mdparagraph:P340:end

\reserve{7\baselineskip}
% mdparagraph:P341:begin
\subsection*{157.5}
% mdparagraph:P341:end

% source-note: FB-VIII-DANGERS-157-N5
% mdparagraph:P342:begin
See 58. 1.
% mdparagraph:P342:end

\reserve{7\baselineskip}
% mdparagraph:P343:begin
\subsection*{157.6}
% mdparagraph:P343:end

% source-note: FB-VIII-DANGERS-157-N6
% mdparagraph:P344:begin
Redundant. What is dūn in Mod. Eng. (24)? Whence is the adverb down derived?
% mdparagraph:P344:end

\reserve{7\baselineskip}
% mdparagraph:P345:begin
\subsection*{157.7}
% mdparagraph:P345:end

% source-note: FB-VIII-DANGERS-157-N7
% mdparagraph:P346:begin
Note the poetical term.
% mdparagraph:P346:end

\reserve{7\baselineskip}
% mdparagraph:P347:begin
\subsection*{158.1}
% mdparagraph:P347:end

% source-note: FB-VIII-DANGERS-158-N1
% mdparagraph:P348:begin
This suggests Seneca (Œdipus, Act 1, v. 11):— Imperia sic excelsa Fortunae objacent.
% mdparagraph:P348:end

\reserve{7\baselineskip}
% mdparagraph:P349:begin
\subsection*{158.2}
% mdparagraph:P349:end

% source-note: FB-VIII-DANGERS-158-N2
% mdparagraph:P350:begin
Pleonastic.
% mdparagraph:P350:end

\reserve{7\baselineskip}
% mdparagraph:P351:begin
\subsection*{158.3}
% mdparagraph:P351:end

% source-note: FB-VIII-DANGERS-158-N3
% mdparagraph:P352:begin
Genitive, dependent on welan.
% mdparagraph:P352:end

\reserve{7\baselineskip}
% mdparagraph:P353:begin
\subsection*{158.4}
% mdparagraph:P353:end

% source-note: FB-VIII-DANGERS-158-N4
% mdparagraph:P354:begin
Cf. Mod. Eng. come to naught.
% mdparagraph:P354:end

\reserve{7\baselineskip}
% mdparagraph:P355:begin
\section*{11. Editorial and grammatical notes}
% mdparagraph:P355:end

% mdparagraph:P356:begin
\textbf{C01 — wæstm.} Cook 26.6 prints wǣstm. The selected corrected form is wæstm, with short æ, as in the earlier noun table. Preserve the original prompt while using the corrected course form. Evidence: Baker’s \emph{Old English Aerobics} glossary, s.v. wæstm; Stenroos (2024), printed 118, Table 1, as checked in the earlier lexical audit.
% mdparagraph:P356:end

% mdparagraph:P357:begin
\textbf{C04 — wund.} Cook 27.4 prints “wund, 43”. The wound/\allowbreak{}sore noun is feminine and takes the selected strong-feminine pattern (§51.b, printed 35 /\allowbreak{} PDF 51), whereas §43 presents masculine nouns. Cook’s glossary adjective wund, “wounded”, is a separate word. Bosworth–Toller’s main volume, printed 1281a, identifies the feminine noun; its Supplement, printed 570b, quotes his wunda in Luke 16.21 and corroborates the accusative plural. The original hint remains printed; no explanation of how the error arose is inferred.
% mdparagraph:P357:end

% mdparagraph:P358:begin
\textbf{C05 — the pointer in note 157.1.} The original says “See p. 156, n. 12.” The parallel explanation “the rest of the forest” is in 156.11, while 156.12 concerns “crag”. This is an apparent reference mismatch; the proposed connection with 156.11 is an inference, not Cook’s stated intention. Independently, ðā ōðre eorðan in the passage means “the rest of the earth”.
% mdparagraph:P358:end

% mdparagraph:P359:begin
\textbf{Quantities and lexical choice.} Cook’s glossary prints dūst and gōd-wębb with long ū and ō; the compound divider is omitted in the course headword gōdwębb. At printed 303 /\allowbreak{} one-based PDF 319, līeg is “thunderbolt; levin”. Accordingly līegum in the reading denotes thunderbolts or lightning. Geneahhe retains the reading’s unmarked spelling; it is not silently replaced by the glossary spelling genēahhe.
% mdparagraph:P359:end

% mdparagraph:P360:begin
\textsuperscript{1}: Cook, \emph{First Book in Old English}, second edition, §§142–147, printed 81–87 (one-based PDF 97–103): selected prefixes, noun and adjective suffixes, composition and the limits of mechanical translation. Lexical meanings follow the selected glossary entries.
% mdparagraph:P360:end

% mdparagraph:P361:begin
\textsuperscript{2}: Cook §70, printed 44 (PDF 60), explicitly compares swīð/\allowbreak{}swīðe and sōð/\allowbreak{}sōðlice. The citation adjectives remain supplied bases.
% mdparagraph:P361:end

% mdparagraph:P362:begin
\textsuperscript{3}: Cook §181, printed 100 (PDF 116), on adjectives as nouns; the personal expression with rihtwīs was taught in Unit 11. Lexical gōd, goods, is separately identified in Cook’s vocabulary, printed 296 (PDF 312), and follows the neuter table.
% mdparagraph:P362:end

% mdparagraph:P363:begin
\textsuperscript{4}: Cook §§43–54, printed 27–38 (PDF 43–54), with the checked glossary entries. The selected reduced wuldr- stem follows §§43.4 and 47.6; reference wuldru follows the disyllabic-neuter route in §47.1. Cook §53.2 gives weak-neuter ēage and genitive-plural ēagena/\allowbreak{}ēagna. These rule-based displays do not claim separate corpus attestations for every reference cell.
% mdparagraph:P363:end

% mdparagraph:P364:begin
\textsuperscript{5}: Cook’s selected weak and strong systems, printed 54–74 (PDF 70–90), with each verb’s glossary entry; present consonant treatment and vowel alternations, printed 24–25 (PDF 40–41). Weak-I gemētt and the selected participle forms follow printed 61–62, 68 and 72. Optional ge- in a participle is discussed at printed 82. Government follows the selected lexical entries and §§156, 164, 172, 184 and 196. The locally supplied speech phrase is supported by Cook printed 153 and 159; the lifetime expression by printed 159.
% mdparagraph:P364:end

% mdparagraph:P365:begin
\textsuperscript{6}: Cook, \emph{First Book in Old English}, second edition, “VIII. The Dangers of Greatness”, printed 156.1–158.4, one-based PDF 172–174. Linguistic text is exact; paragraph joining and removal of source cues are declared layout normalization. The selected edition’s attribution and its qualification are reproduced above. Cook, \emph{Exercises in Old English}, Group VIII, Exercises 26–27, supplies the unchanged English prompts.
% mdparagraph:P365:end

\endgroup
\end{document}
